\documentclass[a4paper,extrafontsizes, 9pt]{memoir}

\usepackage{amsmath, amssymb, amsfonts, amsthm}
\usepackage{multicol}
\usepackage{multirow}
\usepackage{lipsum}
\usepackage{mathtools}
\usepackage{pgfplots}
\usepackage{soul}
\usepackage{hyperref}
\usepackage{enumitem}
\usepackage{fancyhdr}

\pagestyle{empty}

\setlrmarginsandblock{1cm}{1cm}{*}
\setulmarginsandblock{1cm}{1cm}{*}
\checkandfixthelayout

\DeclareMathOperator{\Cov}{Cov}
\DeclareMathOperator{\Var}{Var}

\let\bf\textbf{}

\setlength{\parindent}{0pt}
\renewcommand{\arraystretch}{1.5}

\renewcommand{\thesection}{\arabic{section}}

\newtheorem{kalstok}{Teorema}[section]
\renewcommand{\thekalstok}{\thesection.\arabic{kalstok}}

\theoremstyle{definition}
\newtheorem{definisi}[kalstok]{Definisi}
\newtheorem{lemma}[kalstok]{Lema}
\newtheorem{proposisi}[kalstok]{Proposisi}
\newtheorem{akibat}[kalstok]{Akibat}
\newtheorem{contoh}[kalstok]{Contoh}

\theoremstyle{plain}
\newtheorem{teorema}[kalstok]{Teorema}

\theoremstyle{remark}
\newtheorem{catatan}[kalstok]{Catatan}

\newcommand{\textbetweenrules}[2][.4pt]{%
  \par\vspace{\topsep}
  \noindent\makebox[\textwidth]{%
    \sbox0{#2}%
    \dimen0=.5\dimexpr\ht0+#1\relax
    \dimen2=-.5\dimexpr\ht0-#1\relax
    \leaders\hrule height \dimen0 depth \dimen2\hfill
    \quad #2\quad
    \leaders\hrule height \dimen0 depth \dimen2\hfill
  }\par\nopagebreak\vspace{\topsep}
}

\begin{document}
\footnotesize \textbf{{Cheatsheet -- Kalkulus Stokastik}} \hfill \textit{Teosofi Hidayah Agung/5002221132}
\begin{multicols}{2}
  \section{\small Teori Ukuran}

  \begin{definisi}[Aljabar Himpunan]
    Misalkan $X$ adalah suatu himpunan tak kosong. Aljabar dari himpunan $X$ adalah suatu koleksi $\mathcal{A}$ dari himpunan-himpunan bagian $X$ yang memenuhi:
    \begin{itemize}[leftmargin=*, itemsep=0.5mm]
      \item (A1) $\mathcal{A}$ tidak kosong.
      \item (A2) Jika $E\in\mathcal{A}$, maka komplemennya $E^c:=X\setminus E$ juga elemen $\mathcal{A}$.
      \item (A3) Jika $E,F\in\mathcal{A}$, maka gabungan $E\cup F\in\mathcal{A}$.
    \end{itemize}
  \end{definisi}

  \begin{lemma}
    Misal $X$ himpunan tak kosong. Jika $\mathcal{A}$ aljabar pada $X$ maka berlaku:
    \begin{itemize}[leftmargin=*, itemsep=0.5mm]
      \item (A1$^*$) $X\in\mathcal{A}$.
      \item (A1$^{\#}$) $\varnothing\in\mathcal{A}$.
      \item (A3$^*$) Jika $E,F\in\mathcal{A}$, maka $E\cap F\in\mathcal{A}$.
      \item (A4) Jika $m\in\mathbb{N}$ dan $A_1,A_2,\ldots,A_m\in\mathcal{A}$ maka $\displaystyle\bigcup_{n=1}^{m}A_n\in\mathcal{A}$.
      \item (A5) Jika $m\in\mathbb{N}$ dan $A_1,A_2,\ldots,A_m\in\mathcal{A}$ maka $\displaystyle\bigcap_{n=1}^{m}A_n\in\mathcal{A}$.
    \end{itemize}
  \end{lemma}

  \begin{definisi}[Sigma-Aljabar]
    Misalkan $X$ adalah suatu himpunan tak kosong. Sigma-aljabar dari himpunan $X$ adalah koleksi $\mathcal{A}$ dari himpunan-himpunan bagian $X$ yang memenuhi:
    \begin{itemize}[leftmargin=*, itemsep=0.5mm]
      \item (S1) $\mathcal{A}$ tidak kosong.
      \item (S2) Jika $E\in\mathcal{A}$, maka komplemennya $E^c:=X\setminus E$ juga elemen $\mathcal{A}$.
      \item (S3) Jika $\{E_n\}_{n\in\mathbb{N}}$ barisan himpunan di $\mathcal{A}$, maka $\displaystyle\bigcup_{n\in\mathbb{N}}E_n\in\mathcal{A}$.
    \end{itemize}
    Pasangan $(X,\mathcal{A})$ disebut ruang terukur, dan setiap himpunan dalam $\mathcal{A}$ disebut himpunan terukur.
  \end{definisi}

  \begin{teorema}
    Misalkan $X$ adalah himpunan tak kosong dan $\mathcal{A}$ sigma-aljabar pada $X$. Maka berlaku:
    \begin{itemize}[leftmargin=*, itemsep=0.5mm]
      \item (S1) $\mathcal{A}$ tidak kosong.
      \item (S2) Jika $E\in\mathcal{A}$, maka $E^c\in\mathcal{A}$.
      \item (S3$^*$) Jika $\{E_n\}_{n\in\mathbb{N}}$ barisan himpunan di $\mathcal{A}$, maka $\displaystyle\bigcap_{n\in\mathbb{N}}E_n\in\mathcal{A}$.
    \end{itemize}
  \end{teorema}

  \begin{teorema}[Dekomposisi Atomik]
    Jika $X$ himpunan tak kosong berhingga dan $\mathcal{F}$ suatu sigma-aljabar pada $X$, maka:
    \begin{itemize}[leftmargin=*, itemsep=0.5mm]
      \item (a) Terdapat elemen atom $E$ dengan sifat: (i) $E\in\mathcal{F}$, (ii) $E\neq\varnothing$, dan (iii) jika $B\in\mathcal{F}$ dan $B\subseteq E$, maka $B=\varnothing$ atau $B=E$.
      \item (b) Atom-atom tersebut saling asing.
      \item (c) Setiap $A\in\mathcal{F}$ dapat dituliskan secara unik sebagai gabungan beberapa atom.
    \end{itemize}
  \end{teorema}

  \begin{akibat}
    $\mathcal{F}$ berkorespondensi satu-satu dengan suatu partisi dari $X$.
  \end{akibat}

  \begin{definisi}[Sigma-Aljabar Borel]
    Jika $X$ ruang topologi, maka sigma-aljabar yang dibangkitkan oleh koleksi semua himpunan terbuka $\mathcal{T}$ ditulis $\sigma(\mathcal{T})$.
    Sigma-aljabar ini disebut sigma-aljabar Borel dan dilambangkan $\mathcal{B}(X)$. Setiap anggota $\mathcal{B}(X)$ disebut himpunan Borel.
  \end{definisi}

  \begin{definisi}[Ukuran Positif]
    Misalkan $(X,\mathcal{A})$ ruang terukur. Ukuran positif pada ruang ini adalah fungsi
    \[\mu:\mathcal{A}\to[0,\infty]\]
    yang memenuhi:
    \begin{enumerate}[leftmargin=*, itemsep=0.5mm]
      \item $\mu(\varnothing)=0$.
      \item Aditif terhitung: untuk barisan himpunan saling asing $E_n\in\mathcal{A}$,
            \[\mu\!\left(\bigcup_{n=1}^{\infty}E_n\right)=\sum_{n=1}^{\infty}\mu(E_n).\]
    \end{enumerate}
    Himpunan $(X,\mathcal{A},\mu)$ disebut ruang ukuran.
  \end{definisi}

  \begin{teorema}[Sifat Ukuran]
    Jika $(X,\mathcal{A},\mu)$ ruang ukuran, maka:
    \begin{itemize}[leftmargin=*, itemsep=0.5mm]
      \item (i) Additif berhingga: jika $E_1,E_2,\ldots,E_m\in\mathcal{A}$ saling asing,
            \[\mu\!\left(\bigcup_{n=1}^{m}E_n\right)=\sum_{n=1}^{m}\mu(E_n).\]
      \item (ii) Monotonisitas: jika $E,F\in\mathcal{A}$ dan $E\subseteq F$, maka $\mu(E)\leq\mu(F)$.
      \item (iii) Pengurangan: untuk $E,F\in\mathcal{A}$ dengan $E\subseteq F$ dan $\mu(E)<\infty$,
            \[\mu(F\setminus E)=\mu(F)-\mu(E).\]
      \item (iv) Subadditivitas: untuk $A,B\in\mathcal{A}$ (dengan salah satu ukuran berhingga),
            \[\mu(A\cup B)=\mu(A)+\mu(B)-\mu(A\cap B).\]
    \end{itemize}
  \end{teorema}

  \begin{teorema}[Ketaksamaan Boole]
    Misalkan $(X,\mathcal{A},\mu)$ ruang ukuran. Untuk setiap barisan $E_1,E_2,\ldots\in\mathcal{A}$ berlaku
    \[\mu\!\left(\bigcup_{n=1}^{\infty}E_n\right)\leq\sum_{n=1}^{\infty}\mu(E_n).\]
  \end{teorema}

  \begin{definisi}[Fungsi Terukur]
    Misalkan $(X,\mathcal{A})$ dan $(Y,\mathcal{B})$ adalah ruang terukur.
    Sebuah fungsi $f:X\to Y$ disebut terukur jika untuk setiap $B\in\mathcal{B}$, prapetaannya
    \[f^{-1}(B)=\{x\in X\mid f(x)\in B\}\]
    berada dalam $\mathcal{A}$.
  \end{definisi}

  \begin{catatan}
    Selain notasi $f^{-1}(B)$, sering juga dipakai notasi
    \[\{f\in B\}=\{x\in X\mid f(x)\in B\}.\]
  \end{catatan}

  \begin{teorema}
    Komposisi dari dua fungsi terukur juga menghasilkan fungsi terukur.
  \end{teorema}

  \begin{teorema}
    Misalkan $(X,\mathcal{A})$ dan $(Y,\mathcal{B})$ ruang terukur, dengan $\mathcal{H}$ pembangkit sigma-aljabar $\mathcal{B}$, yaitu $\sigma(\mathcal{H})=\mathcal{B}$.
    Sebuah fungsi $f:X\to Y$ terukur jika, dan hanya jika, untuk setiap $V\in\mathcal{H}$ berlaku
    \[f^{-1}(V)\in\mathcal{A}.\]
  \end{teorema}

  \begin{teorema}
    Misalkan $(X,\mathcal{A})$ ruang terukur. Suatu fungsi $f:X\to\mathbb{R}$ terukur jika dan hanya jika
    \[\{f>c\}=f^{-1}((c,\infty))\in\mathcal{A}\]
    untuk setiap $c\in\mathbb{R}$.
  \end{teorema}

  \begin{contoh}
    Misalkan $(\mathbb{R},\mathcal{B}(\mathbb{R}))$ ruang terukur dengan $\mathcal{B}(\mathbb{R})$ sigma-aljabar Borel pada $\mathbb{R}$.
    Definisikan fungsi
    \[f:\mathbb{R}\to\mathbb{R},\qquad f(x)=x^2.\]
    Untuk setiap interval terbuka $(a,\infty)$ berlaku
    \[f^{-1}((a,\infty))=\{x\in\mathbb{R}\mid x^2>a\}.\]
    Jika $a\ge 0$, maka
    \[f^{-1}((a,\infty))=(-\infty,-\sqrt{a})\cup(\sqrt{a},\infty),\]
    yang merupakan gabungan dua interval terbuka, sehingga termasuk dalam $\mathcal{B}(\mathbb{R})$.
    Karena prapeta setiap himpunan terbuka adalah himpunan Borel, maka $f$ fungsi terukur.
  \end{contoh}

  \begin{catatan}
    Dalam pernyataan teorema sebelumnya, interval $(c,\infty)$ dapat diganti dengan $[c,\infty)$, $(-\infty,c)$,
    atau variasi ekuivalen lainnya, tanpa mengubah esensi.
  \end{catatan}

  \begin{teorema}
    Misalkan $(X,\mathcal{A})$ ruang terukur dan $\{f_n\}$ barisan fungsi terukur $f_n:X\to\overline{\mathbb{R}}$.
    Jika untuk setiap $x\in X$ limit
    \[f(x)=\lim_{n\to\infty}f_n(x)\]
    ada (hingga atau tak hingga), maka fungsi $f$ adalah terukur.
  \end{teorema}

  \begin{definisi}
    Misalkan $S\subseteq X$. Fungsi $\mathbf{1}_S:X\to\mathbb{R}$ yang didefinisikan oleh
    \[\mathbf{1}_S(x)=\begin{cases}
        1, & \text{jika } x\in S,   \\
        0, & \text{jika } x\notin S
      \end{cases}\]
    disebut fungsi indikator (fungsi karakteristik) dari himpunan $S$.
  \end{definisi}

  \begin{catatan}
    Fungsi indikator digunakan untuk merepresentasikan keberadaan titik dalam suatu himpunan.
    Nilainya $1$ jika titik berada dalam himpunan dan $0$ jika berada di luar himpunan.
    Selain itu, $\mathbf{1}_S$ terukur jika dan hanya jika $S$ adalah himpunan terukur.
  \end{catatan}

  \begin{definisi}
    Suatu fungsi $\varphi:X\to\mathbb{R}$ disebut fungsi sederhana jika fungsi tersebut hanya mengambil
    sejumlah hingga nilai. Dengan fungsi indikator, setiap fungsi sederhana dapat direpresentasikan sebagai
    \[\varphi(x)=\sum_{k=1}^{n}a_k\mathbf{1}_{E_k}(x).\]
  \end{definisi}

  \begin{definisi}
    Misalkan $(X,\mathcal{A},\mu)$ ruang ukuran. Integral Lebesgue pada $X$ untuk fungsi sederhana terukur
    bernilai di $\mathbb{R}^{+}$ didefinisikan sebagai
    \[\int_X\varphi\,d\mu=\sum_{k=1}^{n}a_k\mu(E_k),\]
    dengan
    \[\varphi=\sum_{k=1}^{n}a_k\mathbf{1}_{E_k}.\]
  \end{definisi}

  \begin{teorema}[Teorema Aproksimasi]
    Misalkan $(X,\mathcal{A},\mu)$ ruang ukuran. Jika $f:X\to[0,\infty]$ fungsi terukur,
    maka terdapat barisan fungsi sederhana terukur tak negatif $\{\varphi_n\}_{n=1}^{\infty}$ sehingga
    \[\varphi_n(x)\nearrow f(x)\quad\text{untuk setiap }x\in X.\]
    Artinya, untuk setiap $x\in X$ berlaku
    \[0\le\varphi_1(x)\le\varphi_2(x)\le\cdots\quad\text{dan}\quad\lim_{n\to\infty}\varphi_n(x)=f(x).\]
    Jika fungsi $f$ terbatas, maka barisan $\{\varphi_n\}$ dapat dipilih sehingga konvergen seragam ke $f$ pada $X$.
  \end{teorema}

  \begin{definisi}
    Misalkan $(X,\mathcal{A},\mu)$ ruang ukuran dan $f:X\to\mathbb{R}$ fungsi terukur.
    Jika $f$ tidak selalu bernilai tak negatif, maka integral Lebesgue dari $f$ didefinisikan sebagai
    \[\int_X f\,d\mu=\int_X f^{+}\,d\mu-\int_X f^{-}\,d\mu,\]
    dengan
    \[f^{+}(x)=\max\{f(x),0\},\qquad f^{-}(x)=\max\{-f(x),0\}.\]
    Definisi ini berlaku dengan syarat bahwa kedua integral
    \[\int_X f^{+}\,d\mu\quad\text{dan}\quad\int_X f^{-}\,d\mu\]
    tidak bernilai $+\infty$ secara bersamaan.
  \end{definisi}

  \begin{definisi}
    Misalkan $(X,\mathcal{A},\mu)$ ruang ukuran. Suatu himpunan terukur $E\subseteq X$ dikatakan
    memiliki $\mu$-ukuran nol (himpunan nol) jika
    \[\mu(E)=0.\]
  \end{definisi}

  \begin{catatan}
    Pada ruang Euclid $\mathbb{R}^{n}$ dengan ukuran Lebesgue, himpunan yang terdiri dari satu titik memiliki ukuran nol.
    Garis atau kurva dalam $\mathbb{R}^{n}$ untuk $n\ge 2$ juga berukuran nol, dan setiap himpunan terhitung memiliki
    ukuran Lebesgue nol. Secara khusus, fungsi $f$ dikatakan nol hampir di mana-mana jika
    \[\{x\in X:f(x)\neq 0\}\]
    memiliki ukuran nol.
  \end{catatan}

  \begin{teorema}
    Misalkan $(X,\mathcal{A},\mu)$ ruang ukuran dan $f:X\to[0,\infty]$ fungsi terukur.
    Fungsi $f$ bernilai nol hampir di mana-mana jika dan hanya jika
    \[\int_X f\,d\mu=0.\]
  \end{teorema}

  \begin{teorema}[Teorema Konvergensi Monoton]
    Misalkan $(X,\mathcal{A},\mu)$ ruang ukuran. Jika $\{f_n\}$ barisan fungsi terukur tak negatif yang memenuhi
    \[f_1\le f_2\le\cdots\quad\text{dan}\quad f_n(x)\nearrow f(x)\]
    untuk setiap $x\in X$, maka berlaku
    \[\int_X f\,d\mu=\int_X\left(\lim_{n\to\infty}f_n\right)d\mu=\lim_{n\to\infty}\int_X f_n\,d\mu.\]
  \end{teorema}

  \begin{teorema}[Teorema Beppo Levi]
    Jika $\{f_n\}_{n=1}^{\infty}$ barisan fungsi terukur tak negatif pada ruang ukuran $(X,\mathcal{A},\mu)$,
    maka berlaku
    \[\int_X\sum_{n=1}^{\infty}f_n\,d\mu=\sum_{n=1}^{\infty}\int_X f_n\,d\mu.\]
  \end{teorema}

  \begin{definisi}
    Suatu fungsi $f:X\to\mathbb{R}$ dikatakan terintegralkan (Lebesgue integrable) jika fungsi tersebut
    terukur dan memenuhi
    \[\int_X |f|\,d\mu<\infty.\]
  \end{definisi}

  \begin{teorema}[Teorema Konvergensi Dominasi]
    Misalkan $(X,\mathcal{A},\mu)$ ruang ukuran. Andaikan $\{f_n\}$ barisan fungsi terukur
    $f_n:X\to\mathbb{R}$ yang konvergen titik demi titik ke fungsi $f$.
    Jika terdapat fungsi terintegralkan $g:X\to[0,\infty]$ sehingga
    \[|f_n(x)|\le g(x)\quad\text{untuk setiap }n\text{ dan setiap }x\in X,\]
    maka fungsi $f_n$ dan $f$ terintegralkan dan berlaku
    \[\lim_{n\to\infty}\int_X |f_n-f|\,d\mu=0.\]
  \end{teorema}

  \begin{teorema}[Generalisasi Teorema Beppo Levi]
    Misalkan $(X,\mathcal{A},\mu)$ ruang ukuran. Jika $\{f_n\}$ barisan fungsi terukur
    $f_n:X\to\mathbb{R}$ yang memenuhi
    \[\sum_{n=1}^{\infty}\int_X |f_n|\,d\mu<\infty,\]
    maka deret $\sum_{n=1}^{\infty}f_n(x)$ konvergen hampir di mana-mana terhadap suatu fungsi
    $f$ yang terintegralkan, dan berlaku
    \[\int_X\sum_{n=1}^{\infty}f_n\,d\mu=\sum_{n=1}^{\infty}\int_X f_n\,d\mu.\]
  \end{teorema}

  \begin{definisi}
    Suatu koleksi himpunan $\mathcal{P}$ disebut $\pi$-sistem apabila tertutup terhadap irisan berhingga,
    yaitu untuk setiap $A,B\in\mathcal{P}$ berlaku
    \[A\cap B\in\mathcal{P}.\]
    Sedangkan koleksi himpunan $\mathcal{L}$ disebut $\lambda$-sistem (Dynkin system) apabila:
    \begin{enumerate}[leftmargin=*, itemsep=0.5mm]
      \item $\Omega\in\mathcal{L}$.
      \item Jika $A\in\mathcal{L}$, maka $A^c\in\mathcal{L}$.
      \item Jika $\{A_n\}_{n=1}^{\infty}\subseteq\mathcal{L}$ saling asing, maka
            \[\bigcup_{n=1}^{\infty}A_n\in\mathcal{L}.\]
    \end{enumerate}
  \end{definisi}

  \begin{teorema}[Teorema $\pi$-$\lambda$ (Dynkin)]
    Misalkan $\mathcal{P}$ suatu $\pi$-sistem dan $\mathcal{L}$ suatu $\lambda$-sistem dengan
    $\mathcal{P}\subseteq\mathcal{L}$. Maka berlaku
    \[\sigma(\mathcal{P})\subseteq\mathcal{L}.\]
  \end{teorema}

  \begin{teorema}[Teorema Tonelli]
    Misalkan $(X,\mathcal{A},\mu)$ dan $(Y,\mathcal{B},\nu)$ adalah dua ruang ukuran
    $\sigma$-hingga. Jika $f:X\times Y\to[0,\infty]$ fungsi terukur, maka:
    \begin{enumerate}[leftmargin=*, itemsep=0.5mm]
      \item Untuk setiap $x\in X$, fungsi $y\mapsto f(x,y)$ terukur terhadap $\nu$, dan untuk setiap
            $y\in Y$, fungsi $x\mapsto f(x,y)$ terukur terhadap $\mu$.
      \item Fungsi berikut terdefinisi (mungkin tak hingga):
            \[x\mapsto\int_Y f(x,y)\,d\nu(y),\qquad y\mapsto\int_X f(x,y)\,d\mu(x).\]
      \item Berlaku persamaan
            \[\int_{X\times Y} f(x,y)\,d(\mu\otimes\nu)(x,y)
              =\int_X\left(\int_Y f(x,y)\,d\nu(y)\right)d\mu(x)
              =\int_Y\left(\int_X f(x,y)\,d\mu(x)\right)d\nu(y).\]
    \end{enumerate}
  \end{teorema}

  \begin{teorema}
    Misalkan $(X,\mathcal{A},\mu)$ ruang ukuran, dan misalkan $g:X\to[0,\infty]$ fungsi terukur.
    Jika didefinisikan fungsi $\nu:\mathcal{A}\to[0,\infty]$ dengan
    \[\nu(E)=\int_E g\,d\mu,\qquad E\in\mathcal{A},\]
    maka $\nu$ merupakan suatu ukuran pada $(X,\mathcal{A})$. Selanjutnya, untuk setiap fungsi terukur
    $f$ berlaku
    \[\int_X f\,d\nu=\int_X fg\,d\mu,\]
    dengan pengertian bahwa kedua ruas terdefinisi (tidak menghasilkan bentuk tak tentu).
    Penulisan ini sering disingkat sebagai
    \[d\nu=g\,d\mu.\]
  \end{teorema}

  \begin{teorema}[Perubahan Variabel]
    Misalkan $(X,\mathcal{A},\mu)$ ruang ukuran dan $(Y,\mathcal{B})$ ruang terukur.
    Jika $f:Y\to\mathbb{R}$ dan $g:X\to Y$ fungsi terukur, maka berlaku
    \[\int_X (f\circ g)\,d\mu=\int_Y f\,d\nu,\]
    dengan $\nu$ ukuran pada $(Y,\mathcal{B})$ yang diberikan oleh
    \[\nu(B)=\mu\bigl(g^{-1}(B)\bigr),\qquad B\in\mathcal{B}.\]
  \end{teorema}

  \begin{teorema}[Perubahan Variabel Diferensial dalam $\mathbb{R}^{n}$]
    Misalkan $g:X\to Y$ suatu difeomorfisme antara himpunan terbuka di $\mathbb{R}^{n}$.
    Jika $A\subseteq X$ himpunan terukur dan $f:Y\to\mathbb{R}$ fungsi terukur, maka berlaku
    \[\int_{g(A)} f(y)\,d\lambda(y)
      =\int_A (f\circ g)(x)\,|\det Dg(x)|\,d\lambda(x).\]
  \end{teorema}

  \begin{teorema}[Teorema Dasar Pertama Kalkulus]
    Misalkan $I\subseteq\mathbb{R}$ suatu interval dan $f:I\to\mathbb{R}$ fungsi yang terintegralkan Lebesgue.
    Jika didefinisikan fungsi $F:I\to\mathbb{R}$ oleh
    \[F(x)=\int_a^x f(t)\,dt,\]
    untuk suatu $a\in I$, maka $F$ bersifat kontinu pada $I$.
    Lebih lanjut, jika $f$ kontinu di titik $x\in I$, maka $F$ terdiferensialkan di titik tersebut dan
    \[F'(x)=f(x).\]
  \end{teorema}

  \begin{teorema}[Ketergantungan Kontinu pada Parameter Integral]
    Misalkan $(X,\mathcal{A},\mu)$ ruang ukuran, $T$ ruang metrik, dan $f:X\times T\to\mathbb{R}$ fungsi
    sehingga untuk setiap $t\in T$, fungsi $f(\cdot,t)$ terukur pada $X$. Definisikan
    \[F(t)=\int_X f(x,t)\,d\mu(x).\]
    Fungsi $F$ kontinu di titik $t_0\in T$ apabila dipenuhi syarat:
    \begin{itemize}[leftmargin=*, itemsep=0.5mm]
      \item (i) Untuk setiap $x\in X$, fungsi $t\mapsto f(x,t)$ kontinu di $t_0$.
      \item (ii) Terdapat fungsi terintegralkan $g:X\to[0,\infty]$ sehingga
            \[|f(x,t)|\le g(x)\quad\text{untuk semua }x\in X\text{ dan }t\text{ di sekitar }t_0.\]
    \end{itemize}
  \end{teorema}

  \begin{definisi}
    Dua bilangan $p$ dan $q$ dalam interval $(1,\infty)$ disebut eksponen konjugat apabila memenuhi
    \[\frac{1}{p}+\frac{1}{q}=1.\]
  \end{definisi}

  \begin{teorema}[Ketaksamaan Holder]
    Jika $f$ dan $g$ fungsi terukur pada ruang ukuran $(X,\mathcal{A},\mu)$,
    dengan $1<p,q<\infty$ dan
    \[\frac{1}{p}+\frac{1}{q}=1,\]
    maka berlaku
    \[\int_X |fg|\,d\mu\le\|f\|_p\,\|g\|_q.\]
  \end{teorema}

  \begin{teorema}[Ketaksamaan Minkowski]
    Misalkan $f$ dan $g$ fungsi terukur pada ruang ukuran $(X,\mathcal{A},\mu)$, dengan $1\le p<\infty$.
    Maka berlaku
    \[\|f+g\|_p\le\|f\|_p+\|g\|_p.\]
  \end{teorema}

  \begin{definisi}
    Suatu ruang ukuran $(X,\mathcal{A},\mu)$ dikatakan memiliki ukuran berhingga jika
    \[\mu(X)<\infty.\]
  \end{definisi}

  \begin{teorema}
    Misalkan $(X,\mathcal{A},\mu)$ ruang ukuran dengan $\mu(X)<\infty$. Jika $1\le r<p<\infty$, maka berlaku
    \[L^p(X)\subseteq L^r(X).\]
    Selain itu, inklusi ini bersifat kontinu: terdapat konstanta $C>0$ sehingga
    \[\|f\|_r\le C\|f\|_p\quad\text{untuk setiap }f\in L^p(X).\]
  \end{teorema}

  \begin{teorema}
    Misalkan $\{f_n\}$ barisan fungsi terukur $f_n:X\to\mathbb{R}$ yang konvergen ke $f$ hampir di semua titik.
    Jika terdapat fungsi $g\in L^p(X)$, dengan $1\le p<\infty$, yang memenuhi
    \[|f_n(x)|\le g(x)\quad\text{untuk hampir semua }x\in X\text{ dan setiap }n,\]
    maka berlaku:
    \begin{enumerate}[leftmargin=*, itemsep=0.5mm]
      \item $f_n\in L^p(X)$ untuk setiap $n$,
      \item $f\in L^p(X)$,
      \item $f_n\to f$ dalam norma $L^p$, yaitu
            \[\lim_{n\to\infty}\|f_n-f\|_p=0.\]
    \end{enumerate}
  \end{teorema}

  \begin{definisi}
    Misalkan $X$ ruang ukur, dan $f:X\to\mathbb{R}$ fungsi terukur.
    Bilangan $M\in[0,\infty]$ disebut batas atas hampir di mana-mana untuk $|f|$ apabila
    \[|f(x)|\le M\quad\text{untuk hampir semua }x\in X.\]
    Infimum dari semua bilangan tersebut disebut supremum esensial dari $|f|$ dan dilambangkan
    \[\|f\|_{\infty}=\inf\{M\in[0,\infty]:|f(x)|\le M\text{ hampir di semua }x\in X\}.\]
  \end{definisi}

  \begin{definisi}
    Ruang $L^{\infty}$ didefinisikan sebagai himpunan semua fungsi terukur $f:X\to\mathbb{R}$ yang memenuhi
    \[\|f\|_{\infty}<\infty.\]
    Norma pada ruang ini diberikan oleh $\|f\|_{\infty}$.
  \end{definisi}

  \begin{catatan}
    Konsep serupa berlaku untuk vektor $\mathbf{a}=(a_1,\ldots,a_n)\in\mathbb{R}^{n}$:
    \[\lim_{p\to\infty}\|\mathbf{a}\|_{p}=\lim_{p\to\infty}\left(|a_1|^p+\cdots+|a_n|^p\right)^{1/p}
      =\max(|a_1|,\ldots,|a_n|)=\|\mathbf{a}\|_{\infty}.\]
  \end{catatan}

  \begin{teorema}[Teorema Riemann-Lebesgue]
    Untuk setiap $f\in L^1(\mathbb{R})$,
    \[\lim_{|\omega|\to\infty}\int_{\mathbb{R}}f(x)\sin(\omega x)\,dx=0.\]
  \end{teorema}

  \section{\small Probabilitas}

  \begin{definisi}
    Misalkan $\mathcal{F}$ suatu sigma-aljabar pada himpunan $\Omega$. Suatu fungsi
    \[\mathbb{P}:\mathcal{F}\to[0,1]\]
    disebut ukuran probabilitas (probability measure) apabila memenuhi:
    \begin{enumerate}[leftmargin=*, itemsep=0.5mm]
      \item Normalisasi: $\mathbb{P}(\Omega)=1$.
      \item Aditif terhitung: jika $\{A_n\}_{n=1}^{\infty}$ barisan himpunan dalam $\mathcal{F}$ yang saling asing
            (yaitu $A_i\cap A_j=\varnothing$ untuk setiap $i\ne j$), maka
            \[\mathbb{P}\!\left(\bigcup_{n=1}^{\infty}A_n\right)=\sum_{n=1}^{\infty}\mathbb{P}(A_n).\]
    \end{enumerate}
    Tripel $(\Omega,\mathcal{F},\mathbb{P})$ disebut ruang probabilitas.
    Setiap himpunan $A\in\mathcal{F}$ disebut kejadian (event). Kejadian $A$ dikatakan terjadi hampir pasti
    (almost surely, a.s.) apabila
    \[\mathbb{P}(A)=1.\]
  \end{definisi}

  \begin{lemma}[Borel-Cantelli I]
    Misalkan $\{A_n\}_{n=1}^{\infty}$ barisan kejadian dalam ruang probabilitas $(\Omega,\mathcal{F},\mathbb{P})$
    yang memenuhi
    \[\displaystyle\sum_{n=1}^{\infty}\mathbb{P}(A_n)<\infty.\]
    Jika
    $\displaystyle B_n=\bigcup_{k=n}^{\infty}A_k,$
    maka
    $\displaystyle \mathbb{P}\!\left(\bigcap_{n=1}^{\infty}B_n\right)=0.$
  \end{lemma}

  \begin{definisi}
    Misalkan $(\Omega,\mathcal{F},\mathbb{P})$ suatu ruang probabilitas.
    Suatu fungsi $\xi:\Omega\to\mathbb{R}$ disebut variabel acak (random variable) apabila $\xi$ terukur,
    yaitu untuk setiap $B\in\mathcal{B}(\mathbb{R})$ berlaku
    \[\{\xi\in B\}\in\mathcal{F}.\]
  \end{definisi}

  \begin{definisi}
    Misalkan $(\Omega,\mathcal{F},\mathbb{P})$ ruang probabilitas, dan $\xi:\Omega\to\mathbb{R}$ variabel acak.
    Sigma-aljabar yang dibangkitkan oleh $\xi$, dilambangkan dengan $\sigma(\xi)$, didefinisikan sebagai
    \[\sigma(\xi):=\sigma\bigl(\{\{\xi\in B\}:B\in\mathcal{B}(\mathbb{R})\}\bigr).\]
    Dengan kata lain, $\sigma(\xi)$ terdiri atas semua himpunan berbentuk $\{\xi\in B\}$,
    dengan $B$ himpunan Borel di $\mathbb{R}$.
  \end{definisi}

  \begin{teorema}
    Misalkan $(\Omega,\mathcal{F},\mathbb{P})$ suatu ruang probabilitas, dan misalkan $\mathcal{H}$
    koleksi himpunan bagian dari $\mathbb{R}$ yang membangkitkan sigma-aljabar Borel, yaitu
    \[\sigma(\mathcal{H})=\mathcal{B}(\mathbb{R}).\]
    Jika $\xi:\Omega\to\mathbb{R}$ variabel acak, maka
    \[\sigma(\xi)=\sigma\bigl(\{\{\xi\in B\}:B\in\mathcal{H}\}\bigr).\]
  \end{teorema}

  \begin{lemma}
    Misalkan $(\Omega,\mathcal{F},\mathbb{P})$ ruang probabilitas. Jika $\xi$ variabel acak dengan
    himpunan nilai $\xi(\Omega)=\{x_1,x_2,\ldots\}$ yang terhitung, maka setiap $A\in\sigma(\xi)$ dapat
    dinyatakan sebagai gabungan terhitung dari himpunan-himpunan $\{\xi=x_n\}$.
    Artinya, terdapat himpunan indeks $I$ sehingga
    \[A=\bigcup_{n\in I}\{\xi=x_n\},\]
    dengan gabungan saling asing.
  \end{lemma}

  \begin{definisi}
    Misalkan $(\Omega,\mathcal{F},\mathbb{P})$ ruang probabilitas, dan
    $\{\xi_i:i\in I\}$ suatu keluarga variabel acak dengan $I$ himpunan indeks.
    Sigma-aljabar yang dihasilkan keluarga tersebut, dilambangkan
    \[\sigma(\{\xi_i:i\in I\}),\]
    didefinisikan sebagai sigma-aljabar terkecil yang memuat semua himpunan berbentuk
    \[\{\xi_i\in B\},\]
    dengan $B\in\mathcal{B}(\mathbb{R})$ dan $i\in I$.
    Secara ekuivalen,
    \[\sigma(\{\xi_i:i\in I\})=\sigma\bigl(\{\{\xi_i\in B\}:B\in\mathcal{B}(\mathbb{R}),\ i\in I\}\bigr).\]
  \end{definisi}

  \begin{lemma}
    Suatu fungsi $f:\mathbb{R}\to\mathbb{R}$ disebut fungsi Borel apabila untuk setiap himpunan Borel
    $B\in\mathcal{B}(\mathbb{R})$ berlaku
    \[f^{-1}(B)\in\mathcal{B}(\mathbb{R}).\]
    Jika $f$ fungsi Borel dan $\xi$ variabel acak, maka komposisi $f(\xi)$ merupakan variabel acak
    yang terukur terhadap $\sigma(\xi)$.
  \end{lemma}

  \begin{lemma}[Doob-Dynkin]
    Misalkan $(\Omega,\mathcal{F},\mathbb{P})$ ruang probabilitas, dan $\xi:\Omega\to\mathbb{R}$
    variabel acak. Jika $\eta:\Omega\to\mathbb{R}$ variabel acak yang $\sigma(\xi)$-terukur,
    maka terdapat fungsi Borel $f:\mathbb{R}\to\mathbb{R}$ sehingga
    \[\eta=f(\xi)\quad\text{hampir di semua titik}.\]
  \end{lemma}

  \begin{definisi}
    Misalkan $(\Omega,\mathcal{F},\mathbb{P})$ ruang probabilitas, dan $\xi:\Omega\to\mathbb{R}$ variabel acak.
    Ukuran probabilitas yang didapatkan dari $\xi$, dilambangkan $\mathbb{P}_{\xi}$,
    didefinisikan pada $\mathcal{B}(\mathbb{R})$ sebagai
    \[\mathbb{P}_{\xi}(B)=\mathbb{P}(\{\xi\in B\}),\qquad B\in\mathcal{B}(\mathbb{R}).\]
    Ukuran probabilitas $\mathbb{P}_{\xi}$ disebut distribusi dari $\xi$.
    Selanjutnya fungsi $F_{\xi}:\mathbb{R}\to[0,1]$ yang didefinisikan oleh
    \[F_{\xi}(x)=\mathbb{P}(\{\xi\le x\})\]
    disebut cumulative distribution function (CDF) dari $\xi$.
  \end{definisi}

  \begin{definisi}
    Misalkan $(\Omega,\mathcal{F},\mathbb{P})$ ruang probabilitas, dan $\xi:\Omega\to\mathbb{R}$ variabel acak.

    \textbf{Distribusi kontinu absolut:} jika terdapat fungsi Borel
    $f_{\xi}:\mathbb{R}\to[0,\infty)$ sehingga untuk setiap $B\in\mathcal{B}(\mathbb{R})$ berlaku
    \[\mathbb{P}(\{\xi\in B\})=\int_B f_{\xi}(x)\,dx,\]
    maka $\xi$ dikatakan memiliki distribusi kontinu absolut, dan $f_{\xi}$ disebut
    probability density function (PDF) dari $\xi$.

    \textbf{Distribusi diskrit:} jika terdapat himpunan terhitung
    $S=\{x_1,x_2,\ldots\}\subset\mathbb{R}$ dengan elemen-elemen saling berbeda sehingga
    untuk setiap $B\in\mathcal{B}(\mathbb{R})$ berlaku
    \[\mathbb{P}(\{\xi\in B\})=\sum_{x\in B\cap S}\mathbb{P}(\{\xi=x\}),\]
    maka $\xi$ dikatakan memiliki distribusi diskrit, dengan probability mass function (PMF)
    \[p_{\xi}(x_i):=\mathbb{P}(\{\xi=x_i\}),\qquad i=1,2,\ldots.\]
  \end{definisi}

  \begin{definisi}
    Misalkan $(\Omega,\mathcal{F},\mathbb{P})$ ruang probabilitas, dan
    $\xi_1,\ldots,\xi_n:\Omega\to\mathbb{R}$ variabel acak.
    Distribusi gabungan dari $\xi_1,\ldots,\xi_n$ didefinisikan sebagai ukuran probabilitas
    $\mathbb{P}_{\xi_1,\ldots,\xi_n}$ pada $\mathbb{R}^n$ yang diberikan oleh
    \[\mathbb{P}_{\xi_1,\ldots,\xi_n}(B)=\mathbb{P}(\{(\xi_1,\ldots,\xi_n)\in B\}),\qquad B\in\mathcal{B}(\mathbb{R}^n).\]

    \textbf{Distribusi kontinu absolut:} jika terdapat fungsi Borel
    $f_{\xi_1,\ldots,\xi_n}:\mathbb{R}^n\to[0,\infty)$ sehingga untuk setiap $B\in\mathcal{B}(\mathbb{R}^n)$
    berlaku
    \[\mathbb{P}(\{(\xi_1,\ldots,\xi_n)\in B\})=\int_B f_{\xi_1,\ldots,\xi_n}(x_1,\ldots,x_n)\,dx_1\cdots dx_n,\]
    maka $f_{\xi_1,\ldots,\xi_n}$ disebut fungsi kerapatan peluang gabungan (PDF bersama).

    \textbf{Distribusi diskrit:} jika terdapat himpunan terhitung $S\subset\mathbb{R}^n$ sehingga
    \[\mathbb{P}(\{(\xi_1,\ldots,\xi_n)\in B\})=\sum_{x\in B\cap S}\mathbb{P}(\{(\xi_1,\ldots,\xi_n)=x\}),\]
    untuk setiap $B\in\mathcal{B}(\mathbb{R}^n)$, maka distribusi tersebut disebut distribusi diskrit.
    Fungsi
    \[p_{\xi_1,\ldots,\xi_n}(x_1,\ldots,x_n):=\mathbb{P}(\{\xi_1=x_1,\ldots,\xi_n=x_n\})\]
    disebut fungsi massa peluang gabungan (PMF bersama), dengan syarat
    \[p_{\xi_1,\ldots,\xi_n}(x_1,\ldots,x_n)\ge0,\qquad
      \sum_{(x_1,\ldots,x_n)\in S}p_{\xi_1,\ldots,\xi_n}(x_1,\ldots,x_n)=1.\]
  \end{definisi}

  \begin{lemma}
    Misalkan $(\Omega,\mathcal{F},\mathbb{P})$ ruang probabilitas, dan
    $\xi_1,\ldots,\xi_n:\Omega\to\mathbb{R}$ variabel acak.
    \begin{enumerate}[leftmargin=*, itemsep=0.5mm]
      \item Jika $\xi_1,\ldots,\xi_n$ berdistribusi kontinu dengan PDF bersama
            $f_{\xi_1,\ldots,\xi_n}$, maka PDF marginal dari $\xi_k$ diberikan oleh
            \[f_{\xi_k}(x_k)=\int_{\mathbb{R}^{n-1}} f_{\xi_1,\ldots,\xi_n}(x_1,\ldots,x_n)
              \,dx_1\cdots dx_{k-1}dx_{k+1}\cdots dx_n.\]
      \item Jika $\xi_1,\ldots,\xi_n$ berdistribusi diskrit dengan PMF bersama
            $p_{\xi_1,\ldots,\xi_n}$ dan $(\xi_1,\ldots,\xi_n)(\Omega)=S$, maka PMF marginal dari $\xi_k$ diberikan oleh
            \[p_{\xi_k}(x_k)=\sum_{\substack{(x_1,\ldots,x_n)\in S\\x_k\ \text{tetap}}}
              p_{\xi_1,\ldots,\xi_n}(x_1,\ldots,x_n).\]
    \end{enumerate}
    Fungsi $f_{\xi_k}$ dan $p_{\xi_k}$ masing-masing disebut distribusi marginal dari variabel acak $\xi_k$.
  \end{lemma}

  \begin{teorema}
    Misalkan $X=(X_1,\ldots,X_k)$ vektor variabel acak diskrit dengan PMF bersama
    $f_X(x_1,\ldots,x_k)$. Misalkan $Y=(Y_1,\ldots,Y_k)$ vektor variabel acak diskrit
    didefinisikan melalui transformasi
    \[Y=u(X),\]
    dengan $u$ suatu fungsi.
    \begin{enumerate}[leftmargin=*, itemsep=0.5mm]
      \item Jika transformasi $u$ bersifat satu-ke-satu, maka PMF bersama dari $Y$ diberikan oleh
            \[f_Y(y_1,\ldots,y_k)=f_X(x_1,\ldots,x_k),\]
            dengan $(x_1,\ldots,x_k)$ solusi unik dari persamaan $y=u(x)$.
      \item Jika transformasi $u$ tidak satu-ke-satu, maka PMF bersama dari $Y$ diberikan oleh
            \[f_Y(y)=\sum_{x\in u^{-1}(y)}f_X(x),\]
            yaitu penjumlahan atas semua nilai $x$ yang dipetakan ke $y$.
    \end{enumerate}
  \end{teorema}

  \begin{teorema}
    Misalkan $X=(X_1,\ldots,X_k)$ vektor variabel acak kontinu dengan PDF bersama
    $f_X(x_1,\ldots,x_k)>0$ pada himpunan $A\subseteq\mathbb{R}^k$.
    Misalkan $Y=(Y_1,\ldots,Y_k)$ vektor variabel acak kontinu yang didefinisikan melalui
    transformasi satu-ke-satu
    \[Y_i=u_i(X_1,\ldots,X_k),\qquad i=1,\ldots,k,\]
    dengan invers $X=w(Y)$. Jika determinan Jacobian dari transformasi invers kontinu dan tidak nol
    pada daerah yang dipertimbangkan, maka PDF bersama dari $Y$ diberikan oleh
    \[f_Y(y_1,\ldots,y_k)=f_X\bigl(w(y_1,\ldots,y_k)\bigr)
      \left|\det\!\left(\frac{\partial(x_1,\ldots,x_k)}{\partial(y_1,\ldots,y_k)}\right)\right|.\]
  \end{teorema}

  \begin{definisi}[Terintegralkan]
    Misalkan $(\Omega,\mathcal{F},\mathbb{P})$ ruang probabilitas.
    Suatu variabel acak $\xi:\Omega\to\mathbb{R}$ disebut terintegralkan (integrable) apabila
    \[\int_{\Omega}|\xi|\,d\mathbb{P}<\infty.\]
    Dalam hal ini, ekspektasi dari $\xi$ didefinisikan sebagai
    \[\mathbb{E}(\xi)=\int_{\Omega}\xi\,d\mathbb{P}.\]
    Himpunan semua variabel acak yang terintegralkan dinotasikan dengan
    \[L^1=\left\{\xi:\int_{\Omega}|\xi|\,d\mathbb{P}<\infty\right\}.\]
    Lebih detail, $L^1$ juga bisa ditulis sebagai $L^1(\Omega)$ atau $L^1(\Omega,\mathcal{F},\mathbb{P})$.
    Untuk setiap $\xi\in L^1(\Omega,\mathcal{F},\mathbb{P})$, didefinisikan norma
    \[\|\xi\|_1=\mathbb{E}(|\xi|)=\int_{\Omega}|\xi|\,d\mathbb{P}.\]
  \end{definisi}

  \begin{lemma}
    Misalkan $(\Omega,\mathcal{F},\mathbb{P})$ ruang probabilitas, dan
    $\xi_1,\xi_2,\ldots,\xi_n$ variabel acak. Jika $h:\mathbb{R}^n\to\mathbb{R}$ fungsi Borel
    sehingga $h(\xi_1,\xi_2,\ldots,\xi_n)$ terintegralkan, maka berlaku
    \[\mathbb{E}\bigl(h(\xi_1,\xi_2,\ldots,\xi_n)\bigr)
      =\int_{\mathbb{R}^n}h(x_1,x_2,\ldots,x_n)\,d\mathbb{P}_{\xi_1,\ldots,\xi_n}(x_1,x_2,\ldots,x_n).\]
  \end{lemma}

  \begin{akibat}
    Misalkan $(\Omega,\mathcal{F},\mathbb{P})$ ruang probabilitas, dan
    $\xi_1,\xi_2,\ldots,\xi_n$ variabel acak. Dinyatakan
    \[\xi=(\xi_1,\xi_2,\ldots,\xi_n).\]
    Jika $\xi$ memiliki distribusi kontinu dengan PDF bersama $f_{\xi}$, dan
    $h:\mathbb{R}^n\to\mathbb{R}$ fungsi Borel sehingga $h(\xi)$ terintegralkan, maka
    \[\mathbb{E}(h(\xi))=\int_{\mathbb{R}^n}h(x_1,x_2,\ldots,x_n)f_{\xi}(x_1,x_2,\ldots,x_n)\,dx_1dx_2\cdots dx_n.\]
    Selanjutnya, jika $\xi$ memiliki distribusi diskrit dengan
    \[\xi(\Omega)=\{y_1,y_2,\ldots\},\]
    dan $y_i\in\mathbb{R}^n$ saling berbeda, maka berlaku
    \[\mathbb{E}(h(\xi))=\sum_i h(y_i)\,\mathbb{P}(\{\xi=y_i\}).\]
  \end{akibat}

  \begin{definisi}[Terintegralkan Kuadrat]
    Misalkan $(\Omega,\mathcal{F},\mathbb{P})$ ruang probabilitas.
    Suatu variabel acak $\xi:\Omega\to\mathbb{R}$ disebut terintegralkan kuadrat
    (square integrable) apabila memenuhi
    \[\int_{\Omega}|\xi|^2\,d\mathbb{P}<\infty.\]
    Himpunan semua variabel acak yang terintegralkan kuadrat dinotasikan dengan
    \[L^2=\left\{\xi:\int_{\Omega}|\xi|^2\,d\mathbb{P}<\infty\right\}.\]
    Lebih detail, $L^2$ juga bisa dituliskan sebagai $L^2(\Omega)$ atau
    $L^2(\Omega,\mathcal{F},\mathbb{P})$. Untuk setiap $\xi\in L^2(\Omega,\mathcal{F},\mathbb{P})$,
    didefinisikan norma
    \[\|\xi\|_2=\left(\int_{\Omega}|\xi|^2\,d\mathbb{P}\right)^{1/2}.\]
    Selanjutnya, varians dari $\xi$ didefinisikan sebagai
    \[\Var(\xi)=\int_{\Omega}(\xi-\mathbb{E}(\xi))^2\,d\mathbb{P}.\]
  \end{definisi}

  \begin{lemma}
    Misalkan $(\Omega,\mathcal{F},\mathbb{P})$ ruang probabilitas. Jika
    $\eta:\Omega\to[0,\infty)$ variabel acak berdistribusi kontinu absolut dan memenuhi
    \[\int_{\Omega}\eta^2\,d\mathbb{P}<\infty,\]
    maka berlaku
    \[\mathbb{E}(\eta^2)=2\int_0^{\infty}t\,\mathbb{P}(\eta>t)\,dt.\]
  \end{lemma}

  \begin{definisi}
    Misalkan $(\Omega,\mathcal{F},\mathbb{P})$ ruang probabilitas, dan $A,B\in\mathcal{F}$.
    Jika $\mathbb{P}(B)>0$, maka probabilitas bersyarat dari $A$ terhadap $B$ didefinisikan sebagai
    \[\mathbb{P}(A\mid B)=\frac{\mathbb{P}(A\cap B)}{\mathbb{P}(B)}.\]
  \end{definisi}

  \begin{lemma}
    Misalkan $(\Omega,\mathcal{F},\mathbb{P})$ ruang probabilitas, dan $A\in\mathcal{F}$.
    Jika $\{B_n\}_{n\ge1}\subseteq\mathcal{F}$ merupakan barisan kejadian saling asing,
    yaitu $B_i\cap B_j=\varnothing$ untuk setiap $i\ne j$, dan memenuhi
    \[\bigcup_{n=1}^{\infty}B_n=\Omega,\qquad \mathbb{P}(B_n)>0\text{ untuk setiap }n,\]
    maka berlaku hukum probabilitas total:
    \[\mathbb{P}(A)=\sum_{n=1}^{\infty}\mathbb{P}(A\mid B_n)\mathbb{P}(B_n).\]
  \end{lemma}

  \begin{definisi}
    Misalkan $(\Omega,\mathcal{F},\mathbb{P})$ ruang probabilitas. Dua kejadian $A$ dan $B$ dalam
    $\mathcal{F}$ disebut independen apabila memenuhi
    \[\mathbb{P}(A\cap B)=\mathbb{P}(A)\mathbb{P}(B).\]
    Secara umum, sejumlah kejadian $A_1,A_2,\ldots,A_n$ dikatakan saling independen apabila
    untuk setiap pilihan indeks $i_1,i_2,\ldots,i_k$ dengan $1\le i_1<i_2<\cdots<i_k\le n$ berlaku
    \[\mathbb{P}(A_{i_1}\cap A_{i_2}\cap\cdots\cap A_{i_k})
      =\mathbb{P}(A_{i_1})\mathbb{P}(A_{i_2})\cdots\mathbb{P}(A_{i_k}).\]
  \end{definisi}

  \begin{lemma}
    Diberikan kejadian $B\in\mathcal{F}$ dengan $\mathbb{P}(B)\neq0$.
    Kejadian $A$ dan $B$ saling independen jika dan hanya jika berlaku
    \[\mathbb{P}(A\mid B)=\mathbb{P}(A).\]
  \end{lemma}

  \begin{definisi}
    Misalkan $(\Omega,\mathcal{F},\mathbb{P})$ ruang probabilitas, dan
    $\xi,\eta:\Omega\to\mathbb{R}$ dua variabel acak.
    Kedua variabel acak tersebut dikatakan independen apabila untuk setiap
    $A,B\in\mathcal{B}(\mathbb{R})$, kejadian
    \[\{\xi\in A\}\quad\text{dan}\quad\{\eta\in B\}\]
    saling independen, yaitu
    \[\mathbb{P}(\{\xi\in A\}\cap\{\eta\in B\})=\mathbb{P}(\{\xi\in A\})\mathbb{P}(\{\eta\in B\}).\]
    Secara umum, variabel acak $\xi_1,\xi_2,\ldots,\xi_n$ dikatakan saling independen apabila
    untuk setiap $B_1,B_2,\ldots,B_n\in\mathcal{B}(\mathbb{R})$ berlaku
    \[\mathbb{P}\!\left(\bigcap_{i=1}^{n}\{\xi_i\in B_i\}\right)=\prod_{i=1}^{n}\mathbb{P}(\{\xi_i\in B_i\}).\]
    Lebih lanjut, keluarga variabel acak $\{\xi_i:i\in I\}$ dikatakan independen apabila
    setiap subkeluarga berhingga $\{\xi_{i_1},\ldots,\xi_{i_k}\}$ dengan
    $i_1,\ldots,i_k\in I$ saling independen.
  \end{definisi}

  \begin{teorema}
    Misalkan $(\Omega,\mathcal{F},\mathbb{P})$ ruang probabilitas, dan
    $\xi,\eta:\Omega\to\mathbb{R}$ dua variabel acak.
    Pernyataan berikut ekuivalen:
    \begin{enumerate}[leftmargin=*, itemsep=0.5mm]
      \item $\xi$ dan $\eta$ saling independen, yaitu untuk setiap $A,B\in\mathcal{B}(\mathbb{R})$ berlaku
            \[\mathbb{P}(\xi\in A,\eta\in B)=\mathbb{P}(\xi\in A)\mathbb{P}(\eta\in B).\]
      \item Distribusi gabungan $\mathbb{P}_{\xi,\eta}$ merupakan hasil kali distribusi marginalnya,
            yaitu
            \[\mathbb{P}_{\xi,\eta}=\mathbb{P}_{\xi}\otimes\mathbb{P}_{\eta}.\]
    \end{enumerate}
    Jika salah satu (dan karenanya kedua) pernyataan di atas berlaku, maka untuk setiap
    fungsi Borel tak-negatif $g$ atau fungsi yang terintegralkan, diperoleh
    \[\int_{\mathbb{R}^2}g(x_1,x_2)\,d\mathbb{P}_{\xi,\eta}(x_1,x_2)
      =\int_{\mathbb{R}}\int_{\mathbb{R}}g(x_1,x_2)\,d\mathbb{P}_{\xi}(x_1)\,d\mathbb{P}_{\eta}(x_2).\]
  \end{teorema}

  \begin{proposisi}
    Misalkan $(\Omega,\mathcal{F},\mathbb{P})$ ruang probabilitas.
    Jika $\xi$ dan $\eta$ dua variabel acak yang terintegralkan dan saling independen,
    serta hasil kali $\xi\eta$ juga terintegralkan, maka berlaku
    \[\mathbb{E}(\xi\eta)=\mathbb{E}(\xi)\mathbb{E}(\eta).\]
    Secara umum, jika $\xi_1,\xi_2,\ldots,\xi_n$ adalah variabel acak yang saling independen
    dan terintegralkan, serta hasil kali $\xi_1\xi_2\cdots\xi_n$ terintegralkan,
    maka berlaku
    \[\mathbb{E}(\xi_1\xi_2\cdots\xi_n)=\prod_{i=1}^{n}\mathbb{E}(\xi_i).\]
  \end{proposisi}

  \begin{teorema}
    Jika $\eta$ dan $\xi$ peubah acak independen pada ruang probabilitas $(\Omega,\mathcal{F},\mathbb{P})$,
    dan $f$ serta $g$ fungsi terukur, maka peubah acak $f(\eta)$ dan $g(\xi)$ juga independen.
  \end{teorema}

  \begin{proposisi}
    Misalkan $\xi$ variabel acak dan $\eta_1,\ldots,\eta_n$ variabel acak.
    Diasumsikan bahwa $\xi$ independen terhadap vektor variabel acak
    \[(\eta_1,\ldots,\eta_n),\]
    yang artinya untuk setiap $B\in\mathcal{B}(\mathbb{R})$ dan setiap $A\in\mathcal{B}(\mathbb{R}^n)$ berlaku
    \[\mathbb{P}(\{\xi\in B,\ g(\eta_1,\ldots,\eta_n)\in A\})
      =\mathbb{P}(\{\xi\in B\})\,\mathbb{P}(\{g(\eta_1,\ldots,\eta_n)\in A\}).\]
    Jika $g:\mathbb{R}^n\to\mathbb{R}^n$ fungsi terukur, maka $\xi$ independen terhadap
    variabel acak vektor $g(\eta_1,\ldots,\eta_n)$.
  \end{proposisi}

  \begin{definisi}
    Misalkan $(\Omega,\mathcal{F},\mathbb{P})$ ruang probabilitas.
    Diberikan dua sigma-aljabar $\mathcal{G}$ dan $\mathcal{H}$ dengan
    \[\mathcal{G},\mathcal{H}\subseteq\mathcal{F}.\]
    Kedua sigma-aljabar tersebut dikatakan independen apabila untuk setiap pasangan kejadian
    $A\in\mathcal{G}$ dan $B\in\mathcal{H}$ berlaku
    \[\mathbb{P}(A\cap B)=\mathbb{P}(A)\mathbb{P}(B).\]
  \end{definisi}

  \begin{lemma}
    Misalkan $(\Omega,\mathcal{F},\mathbb{P})$ ruang probabilitas dan
    $\xi,\eta:\Omega\to\mathbb{R}$ dua variabel acak.
    Maka $\xi$ dan $\eta$ saling independen jika dan hanya jika sigma-aljabar
    $\sigma(\xi)$ dan $\sigma(\eta)$ saling independen.
  \end{lemma}

  \begin{definisi}
    Misalkan $(\Omega,\mathcal{F},\mathbb{P})$ ruang probabilitas, dan
    $\mathcal{G}\subseteq\mathcal{F}$ suatu sigma-aljabar.
    Suatu variabel acak $\xi:\Omega\to\mathbb{R}$ dikatakan independen terhadap sigma-aljabar
    $\mathcal{G}$ apabila sigma-aljabar $\sigma(\xi)$ dan $\mathcal{G}$ saling independen,
    yaitu apabila untuk setiap $A\in\mathcal{B}(\mathbb{R})$ dan setiap $G\in\mathcal{G}$ berlaku
    \[\mathbb{P}(\{\xi\in A\}\cap G)=\mathbb{P}(\xi\in A)\mathbb{P}(G).\]
    Dengan kata lain, nilai variabel acak $\xi$ tidak dipengaruhi oleh informasi yang terkandung
    dalam sigma-aljabar $\mathcal{G}$.
  \end{definisi}

  \section{\small Ekspektasi Bersyarat}

  \begin{definisi}
    Misalkan $(\Omega,\mathcal{F},\mathbb{P})$ adalah suatu ruang probabilitas, dan
    $\xi:\Omega\to\mathbb{R}$ merupakan variabel acak yang terintegralkan, yaitu memenuhi
    \[\int_{\Omega}|\xi|\,d\mathbb{P}<\infty.\]
    Diberikan suatu kejadian $B\in\mathcal{F}$ dengan syarat $\mathbb{P}(B)>0$, ekspektasi
    bersyarat dari $\xi$ terhadap kejadian $B$ didefinisikan sebagai
    \[\mathbb{E}(\xi\mid B)=\frac{1}{\mathbb{P}(B)}\int_B \xi\,d\mathbb{P}.\]
  \end{definisi}

  \begin{definisi}
    Misalkan $(\Omega,\mathcal{F},\mathbb{P})$ adalah suatu ruang probabilitas. Diberikan
    variabel acak $\xi:\Omega\to\mathbb{R}$ yang terintegralkan, dan variabel acak diskrit
    $\eta:\Omega\to\mathbb{R}$ dengan
    \[S:=\eta(\Omega)=\{y_1,y_2,\ldots\},\]
    di mana $y_1,y_2,\ldots$ saling berbeda dan
    \[\mathbb{P}(\{\eta=y_n\})>0,\qquad \text{untuk setiap }n.\]
    Ekspektasi bersyarat dari $\xi$ dengan syarat $\eta$ didefinisikan sebagai variabel acak
    $\mathbb{E}(\xi\mid\eta)$ yang diberikan oleh
    \[\mathbb{E}(\xi\mid\eta)
      =\sum_{n=1}^{\infty}\mathbb{E}(\xi\mid\{\eta=y_n\})\mathbf{1}_{\{\eta=y_n\}}.\]
    Secara ekuivalen, ekspektasi bersyarat tersebut dapat dituliskan sebagai
    \[\mathbb{E}(\xi\mid\eta)(\omega)=\mathbb{E}(\xi\mid\{\eta=\eta(\omega)\}),\qquad \omega\in\Omega.\]
  \end{definisi}

  \begin{lemma}
    Misalkan $(\Omega,\mathcal{F},\mathbb{P})$ adalah suatu ruang probabilitas.
    Jika $\eta$ adalah variabel acak diskrit dan $\xi$ adalah variabel acak yang dapat
    diintegralkan, maka berlaku
    \[\mathbb{E}\bigl(\mathbb{E}(\xi\mid\eta)\bigr)=\mathbb{E}(\xi).\]
  \end{lemma}

  \begin{proposisi}
    Misalkan $(\Omega,\mathcal{F},\mathbb{P})$ adalah suatu ruang probabilitas.
    Jika $\xi$ adalah variabel acak yang terintegralkan dan $\eta$ adalah variabel acak diskrit,
    maka berlaku:
    \begin{enumerate}[leftmargin=*, itemsep=0.5mm]
      \item $\mathbb{E}(\xi\mid\eta)$ adalah $\sigma(\eta)$-terukur.
      \item Untuk setiap $A\in\sigma(\eta)$, berlaku
            \[\int_A \mathbb{E}(\xi\mid\eta)\,d\mathbb{P}=\int_A \xi\,d\mathbb{P}.\]
    \end{enumerate}
  \end{proposisi}

  \begin{definisi}
    Misalkan $(\Omega,\mathcal{F},\mathbb{P})$ adalah suatu ruang probabilitas. Diberikan
    variabel acak $\xi$ yang terintegralkan dan variabel acak $\eta$ dengan nilai yang umum
    (tidak harus diskrit). Ekspektasi bersyarat dari $\xi$ dengan syarat $\eta$, yang
    dinotasikan dengan $\mathbb{E}(\xi\mid\eta)$, didefinisikan sebagai suatu variabel acak
    yang memenuhi sifat-sifat berikut:
    \begin{enumerate}[leftmargin=*, itemsep=0.5mm]
      \item $\mathbb{E}(\xi\mid\eta)$ adalah $\sigma(\eta)$-terukur.
      \item Untuk setiap himpunan $A\in\sigma(\eta)$, berlaku
            \[\int_A \mathbb{E}(\xi\mid\eta)\,d\mathbb{P}=\int_A \xi\,d\mathbb{P}.\]
    \end{enumerate}
  \end{definisi}

  \begin{catatan}
    Selain itu, probabilitas bersyarat dari suatu kejadian $A\in\mathcal{F}$ dengan syarat
    variabel acak $\eta$ dapat didefinisikan melalui ekspektasi bersyarat, yaitu
    \[\mathbb{P}(A\mid\eta)=\mathbb{E}(\mathbf{1}_A\mid\eta).\]
    Perhatikan bahwa $\mathbb{P}(A\mid\eta)$ merupakan variabel acak yang $\sigma(\eta)$-terukur,
    sehingga nilainya hanya bergantung pada informasi yang dikandung oleh $\eta$. Sebagai kasus
    khusus, jika $\eta$ adalah variabel acak diskrit, maka untuk setiap nilai $y_n$ diperoleh
    \[\mathbb{P}(A\mid\eta)(\omega)=\mathbb{P}(A\mid\{\eta=y_n\}),\]
    untuk setiap $\omega$ dengan $\eta(\omega)=y_n$.
  \end{catatan}

  \begin{lemma}
    Misalkan $(\Omega,\mathcal{F},\mathbb{P})$ adalah suatu ruang probabilitas, dan
    $\mathcal{G}\subseteq\mathcal{F}$ adalah suatu sigma-aljabar. Jika $\zeta$ adalah variabel
    acak yang $\mathcal{G}$-terukur dan untuk setiap $B\in\mathcal{G}$ berlaku
    \[\int_B \zeta\,d\mathbb{P}=0,\]
    maka diperoleh
    \[\zeta=0\qquad \text{hampir pasti}.\]
  \end{lemma}

  \begin{lemma}
    Jika $\xi$ dan $\eta$ adalah variabel acak terintegral yang memiliki PDF bersama
    $f_{\xi,\eta}(x,y)$, maka berlaku
    \[\mathbb{E}(\xi\mid\eta)
      =\frac{\int_{\mathbb{R}}x\,f_{\xi,\eta}(x,\eta)\,dx}
      {\int_{\mathbb{R}}f_{\xi,\eta}(x,\eta)\,dx},
      \qquad \text{hampir pasti},\]
    pada himpunan di mana penyebut bernilai positif.
  \end{lemma}

  \begin{proposisi}
    Jika $\sigma(\eta)=\sigma(\eta')$, maka berlaku
    \[\mathbb{E}(\xi\mid\eta)=\mathbb{E}(\xi\mid\eta')\qquad \text{hampir pasti}.\]
  \end{proposisi}


  \begin{definisi}
    Misalkan $\xi$ adalah variabel acak terintegral pada ruang probabilitas
    $(\Omega,\mathcal{F},\mathbb{P})$, yaitu $\mathbb{E}(|\xi|)<\infty$, dan misalkan
    $\mathcal{G}\subseteq\mathcal{F}$ adalah suatu sigma-aljabar. Ekspektasi bersyarat dari
    $\xi$ terhadap $\mathcal{G}$ didefinisikan sebagai suatu variabel acak
    $\mathbb{E}(\xi\mid\mathcal{G})$ yang memenuhi dua kondisi berikut:
    \begin{enumerate}[leftmargin=*, itemsep=0.5mm]
      \item $\mathbb{E}(\xi\mid\mathcal{G})$ terukur terhadap $\mathcal{G}$.
      \item Untuk setiap $A\in\mathcal{G}$ berlaku
            \[\int_A \mathbb{E}(\xi\mid\mathcal{G})\,d\mathbb{P}=\int_A \xi\,d\mathbb{P}.\]
    \end{enumerate}
    Secara intuitif, $\mathbb{E}(\xi\mid\mathcal{G})$ dapat dipandang sebagai "rata-rata terbaik"
    dari $\xi$ yang hanya bergantung pada informasi yang tersedia dalam sigma-aljabar
    $\mathcal{G}$.
  \end{definisi}

  \begin{proposisi}
    Jika $\xi$ adalah variabel acak terintegral pada ruang probabilitas
    $(\Omega,\mathcal{F},\mathbb{P})$, dan $\mathcal{G}\subseteq\mathcal{F}$ adalah suatu
    sigma-aljabar, maka ekspektasi bersyarat $\mathbb{E}(\xi\mid\mathcal{G})$ selalu ada dan
    tunggal hingga hampir pasti. Lebih lanjut, jika $\xi$ dan $\xi'$ adalah dua variabel acak
    terintegral yang memenuhi
    \[\xi=\xi'\qquad \text{hampir pasti},\]
    maka berlaku
    \[\mathbb{E}(\xi\mid\mathcal{G})=\mathbb{E}(\xi'\mid\mathcal{G})\qquad \text{hampir pasti}.\]
  \end{proposisi}

  \begin{teorema}[Teorema Proyeksi pada Ruang Hilbert]
    Misalkan $H$ adalah ruang Hilbert dengan hasil kali dalam $\langle\cdot,\cdot\rangle$, dan
    misalkan $M\subseteq H$ adalah subruang linear tertutup. Maka untuk setiap $X\in H$ terdapat
    elemen unik $\zeta\in M$ sedemikian sehingga
    \[\langle X-\zeta,Z\rangle=0,\qquad \text{untuk setiap }Z\in M.\]
  \end{teorema}

  \begin{akibat}
    Misalkan $(\Omega,\mathcal{F},\mathbb{P})$ adalah ruang probabilitas dan
    $\mathcal{G}\subseteq\mathcal{F}$ adalah suatu sigma-aljabar. Untuk setiap $A\in\mathcal{F}$,
    terdapat suatu variabel acak $\zeta$ yang $\mathcal{G}$-terukur sehingga
    \[\int_B \zeta\,d\mathbb{P}=\mathbb{P}(A\cap B),\qquad \text{untuk setiap }B\in\mathcal{G}.\]
    Variabel acak $\zeta$ tersebut disebut sebagai ekspektasi bersyarat dari $\mathbf{1}_A$
    terhadap $\mathcal{G}$ dan dinotasikan dengan
    \[\zeta=\mathbb{E}(\mathbf{1}_A\mid\mathcal{G}).\]
  \end{akibat}

  \begin{lemma}
    Misalkan $(\Omega,\mathcal{F},\mathbb{P})$ adalah ruang probabilitas dan
    $\mathcal{G}\subseteq\mathcal{F}$ adalah suatu sigma-aljabar. Jika $\{\xi_n\}$ adalah
    barisan variabel acak terintegral tak-negatif yang memenuhi
    \[\xi_n\uparrow \xi\qquad \text{hampir pasti},\]
    dan $\{\zeta_n\}$ adalah barisan variabel acak terintegral yang $\mathcal{G}$-terukur yang
    memenuhi
    \[\int_A \zeta_n\,d\mathbb{P}=\int_A \xi_n\,d\mathbb{P},\qquad \text{untuk setiap }A\in\mathcal{G},\]
    maka berlaku
    \[\zeta_n\uparrow \zeta\qquad \text{hampir pasti}\]
    untuk suatu variabel acak $\zeta$ yang $\mathcal{G}$-terukur dan terintegral.
  \end{lemma}

  \begin{teorema}[Radon-Nikodym]
    Misalkan $(\Omega,\mathcal{F},\mathbb{P})$ adalah ruang probabilitas, dan
    $\mathcal{G}\subseteq\mathcal{F}$ adalah suatu sigma-aljabar. Jika $\xi$ adalah variabel
    acak terintegralkan, maka terdapat suatu variabel acak $\zeta$ yang $\mathcal{G}$-terukur
    sehingga
    \[\int_A \xi\,d\mathbb{P}=\int_A \zeta\,d\mathbb{P},\]
    untuk setiap $A\in\mathcal{G}$.
  \end{teorema}

  \begin{proposisi}
    Ekspektasi bersyarat memiliki beberapa sifat berikut. Misalkan $\xi$ dan $\zeta$ adalah
    variabel acak terintegral pada ruang probabilitas $(\Omega,\mathcal{F},\mathbb{P})$, dan
    $\mathcal{G},\mathcal{H}\subseteq\mathcal{F}$ adalah sigma-aljabar.
    \begin{enumerate}[leftmargin=*, itemsep=0.5mm]
      \item (Linearitas) Untuk setiap $a,b\in\mathbb{R}$ berlaku
            \[\mathbb{E}(a\xi+b\zeta\mid\mathcal{G})
              =a\,\mathbb{E}(\xi\mid\mathcal{G})+b\,\mathbb{E}(\zeta\mid\mathcal{G}).\]
      \item (Konsistensi ekspektasi)
            \[\mathbb{E}\bigl(\mathbb{E}(\xi\mid\mathcal{G})\bigr)=\mathbb{E}(\xi).\]
      \item (Faktorisasi) Jika $\xi$ adalah $\mathcal{G}$-terukur dan $\xi\zeta$ terintegral, maka
            \[\mathbb{E}(\xi\zeta\mid\mathcal{G})=\xi\,\mathbb{E}(\zeta\mid\mathcal{G}).\]
      \item (Independensi) Jika $\xi$ independen terhadap $\mathcal{G}$, maka
            \[\mathbb{E}(\xi\mid\mathcal{G})=\mathbb{E}(\xi).\]
      \item (Sifat menara) Jika $\mathcal{H}\subseteq\mathcal{G}$, maka
            \[\mathbb{E}\bigl(\mathbb{E}(\xi\mid\mathcal{G})\mid\mathcal{H}\bigr)
              =\mathbb{E}(\xi\mid\mathcal{H}).\]
      \item (Monotonisitas) Jika $\xi\ge 0$ hampir pasti, maka
            \[\mathbb{E}(\xi\mid\mathcal{G})\ge 0\qquad \text{hampir pasti}.\]
    \end{enumerate}
  \end{proposisi}

  \begin{teorema}[Ketidaksamaan Jensen]
    Misalkan $\varphi:\mathbb{R}\to\mathbb{R}$ adalah fungsi konveks. Misalkan $\xi$ adalah
    variabel acak terintegral pada ruang probabilitas $(\Omega,\mathcal{F},\mathbb{P})$, dan
    diasumsikan bahwa $\varphi(\xi)$ juga terintegral. Jika $\mathcal{G}\subseteq\mathcal{F}$
    adalah suatu sigma-aljabar, maka berlaku
    \[\varphi\bigl(\mathbb{E}(\xi\mid\mathcal{G})\bigr)\le
      \mathbb{E}\bigl(\varphi(\xi)\mid\mathcal{G}\bigr)\qquad \text{hampir pasti}.\]
  \end{teorema}

  \section{\small Proses Stokastik Waktu Kontinu}

  \begin{definisi}
    Misalkan $(\Omega,\mathcal{F},\mathbb{P})$ adalah suatu ruang probabilitas. Suatu proses
    stokastik didefinisikan sebagai keluarga variabel acak
    \[\{X_t:t\in T\},\]
    yang didefinisikan pada ruang probabilitas tersebut, dengan $T$ merupakan himpunan indeks
    (parameter waktu). Selanjutnya, konsep-konsep berikut didefinisikan:
    \begin{enumerate}[leftmargin=*, itemsep=0.5mm]
      \item Jika $T=\{1,2,3,\ldots\}$, maka proses $\{X_t\}$ disebut sebagai proses stokastik
            waktu diskrit.
      \item Jika $T$ merupakan suatu interval dalam $\mathbb{R}$ (misalnya $T=[0,\infty)$), maka
            proses $\{X_t\}$ disebut sebagai proses stokastik waktu kontinu.
      \item Jika untuk setiap $t\in T$ berlaku
            \[X_t:\Omega\to S,\]
            maka himpunan $S$ disebut sebagai ruang keadaan (state space) dari proses stokastik
            tersebut.
      \item Untuk setiap $\omega\in\Omega$, fungsi
            \[T\ni t\mapsto X_t(\omega)\]
            disebut sebagai lintasan sampel (sample path) dari proses stokastik.
    \end{enumerate}
  \end{definisi}

  \begin{definisi}
    Misalkan $(\Omega,\mathcal{F},\mathbb{P})$ adalah suatu ruang probabilitas dan $T$ adalah
    himpunan indeks (waktu). Suatu keluarga sigma-aljabar $\{\mathcal{F}_t\}_{t\in T}$ pada
    $\Omega$ disebut sebagai filtrasi jika memenuhi
    \[\mathcal{F}_s\subseteq\mathcal{F}_t\subseteq\mathcal{F},\qquad
      \text{untuk setiap }s,t\in T\text{ dengan }s\le t.\]
    Keluarga $\{\mathcal{F}_t\}$ menginterpretasikan informasi yang tersedia hingga waktu $t$,
    sehingga kondisi $\mathcal{F}_s\subseteq\mathcal{F}_t$ menyatakan bahwa informasi tidak
    berkurang seiring bertambahnya waktu.
  \end{definisi}

  \begin{definisi}
    Suatu proses stokastik $\{X_t\}_{t\in T}$ dikatakan sebagai martingale
    (masing-masing submartingale atau supermartingale) terhadap suatu filtrasi
    $\{\mathcal{F}_t\}_{t\in T}$ jika kondisi-kondisi berikut dipenuhi:
    \begin{enumerate}[leftmargin=*, itemsep=0.5mm]
      \item Untuk setiap $t\in T$, variabel acak $X_t$ terintegralkan, yaitu
            \[\mathbb{E}(|X_t|)<\infty.\]
      \item Untuk setiap $t\in T$, variabel acak $X_t$ adalah $\mathcal{F}_t$-terukur. Dengan
            kata lain, proses $\{X_t\}$ bersifat adapted terhadap filtrasi $\{\mathcal{F}_t\}$.
      \item Untuk setiap $s,t\in T$ dengan $s\le t$, berlaku
            \[\mathbb{E}(X_t\mid\mathcal{F}_s)=X_s,\]
            untuk martingale, sedangkan untuk submartingale dan supermartingale masing-masing
            berlaku
            \[\mathbb{E}(X_t\mid\mathcal{F}_s)\ge X_s,\qquad
              \text{atau}\qquad \mathbb{E}(X_t\mid\mathcal{F}_s)\le X_s.\]
    \end{enumerate}
  \end{definisi}

  \begin{definisi}
    Suatu proses stokastik $\{W(t)\}_{t\ge 0}$ yang bernilai dalam $\mathbb{R}$ disebut sebagai
    proses Wiener (atau gerak Brown) apabila memenuhi sifat-sifat berikut:
    \begin{enumerate}[leftmargin=*, itemsep=0.5mm]
      \item Kondisi awal dipenuhi, yaitu hampir pasti berlaku
            \[W(0)=0.\]
      \item Lintasan sampel dari proses $\{W(t)\}_{t\ge 0}$ bersifat kontinu, yaitu untuk hampir
            setiap $\omega$, fungsi $t\mapsto W(t,\omega)$ kontinu pada $[0,\infty)$.
      \item Untuk setiap barisan waktu $0<t_1<t_2<\cdots<t_n$ dan setiap himpunan Borel
            $A_1,A_2,\ldots,A_n\subseteq\mathbb{R}$, distribusi bersama dari
            $(W(t_1),W(t_2),\ldots,W(t_n))$ diberikan oleh
            \begin{align*}
               & \mathbb{P}(W(t_1)\in A_1,\ldots,W(t_n)\in A_n) \\
               & \qquad =\int_{A_1}\cdots\int_{A_n}
              p(t_1,0,x_1)\,p(t_2-t_1,x_1,x_2)\cdots p(t_n-t_{n-1},x_{n-1},x_n)\,
              dx_1\cdots dx_n,
            \end{align*}
            dengan PDF transisi diberikan oleh
            \[p(t,x,y)=\frac{1}{\sqrt{2\pi t}}
              \exp\!\left(-\frac{(y-x)^2}{2t}\right),\]
            untuk setiap $x,y\in\mathbb{R}$ dan $t>0$. Fungsi $p(t,x,y)$ tersebut disebut
            sebagai PDF transisi.
    \end{enumerate}
  \end{definisi}

  \begin{lemma}
    Jika $\{W(t)\}_{t\ge 0}$ adalah proses Wiener, maka untuk setiap $t\ge 0$, variabel acak
    $W(t)$ memiliki PDF
    \[p(t,0,x)=\frac{1}{\sqrt{2\pi t}}\exp\!\left(-\frac{x^2}{2t}\right),\]
    yang artinya berdistribusi normal dengan ekspektasi $0$ dan variansi $t$, yaitu
    \[W(t)\sim\mathcal{N}(0,t).\]
  \end{lemma}

  \begin{lemma}
    Jika $\{W(t)\}_{t\ge 0}$ adalah proses Wiener, maka untuk setiap $s,t\in[0,\infty)$ berlaku
    \[\mathbb{E}(W(s)W(t))=\min\{s,t\}.\]
  \end{lemma}

  \begin{akibat}
    Jika $\{W(t)\}_{t\ge 0}$ adalah proses Wiener, maka untuk setiap $s,t\in[0,\infty)$ berlaku
    \[\mathbb{E}\bigl((W(t)-W(s))^2\bigr)=|t-s|.\]
  \end{akibat}

  \begin{lemma}
    Jika $\{W(t)\}_{t\ge 0}$ adalah proses Wiener, maka fungsi karakteristik dari variabel acak
    $W(t)$ diberikan oleh
    \[\varphi_{W(t)}(\lambda)=\mathbb{E}\bigl(e^{i\lambda W(t)}\bigr)
      =e^{-\lambda^2 t/2},\qquad \lambda\in\mathbb{R}.\]
  \end{lemma}

  \begin{definisi}
    Suatu proses stokastik $\{W(t)\}_{t\ge 0}$ dengan nilai pada $\mathbb{R}^n$ yang dinyatakan
    sebagai
    \[W(t)=\bigl(W^1(t),W^2(t),\ldots,W^n(t)\bigr)\]
    disebut sebagai proses Wiener berdimensi-$n$ apabila setiap komponen
    $\{W^1(t)\}_{t\ge 0},\ldots,\{W^n(t)\}_{t\ge 0}$ merupakan proses Wiener bernilai real dan
    saling bebas, yaitu untuk setiap $t\ge 0$, variabel acak
    $W^1(t),\ldots,W^n(t)$ saling bebas.
  \end{definisi}

  \begin{definisi}
    Suatu proses stokastik $\{\xi(t)\}_{t\in T}$ dikatakan memiliki inkremen yang stasioner
    apabila untuk setiap $s,t\in T$ dan setiap $h$ sedemikian sehingga $s+h,t+h\in T$, berlaku
    \[\xi(t+h)-\xi(s+h)\overset{d}{=}\xi(t)-\xi(s),\]
    di mana $\overset{d}{=}$ menyatakan kesamaan distribusi. Dengan kata lain, distribusi
    probabilitas dari inkremen $\xi(t)-\xi(s)$ hanya bergantung pada selisih waktu $t-s$, dan
    tidak bergantung pada posisi waktu absolut.
  \end{definisi}

  \begin{proposisi}
    Proses stokastik Wiener $\{W(t)\}_{t\ge 0}$ memiliki inkremen yang stasioner, yaitu untuk
    setiap $0\le s<t$, inkremen $W(t)-W(s)$ berdistribusi normal dengan nilai rata-rata $0$ dan
    variansi $t-s$, yaitu
    \[W(t)-W(s)\sim\mathcal{N}(0,t-s).\]
  \end{proposisi}

  \begin{definisi}
    Suatu proses stokastik $\{\xi(t)\}_{t\in T}$ dikatakan memiliki inkremen yang independen
    apabila untuk setiap barisan waktu
    \[t_0<t_1<\cdots<t_n,\qquad t_0,t_1,\ldots,t_n\in T,\]
    variabel acak
    \[\xi(t_1)-\xi(t_0),\ \xi(t_2)-\xi(t_1),\ \ldots,\ \xi(t_n)-\xi(t_{n-1})\]
    saling independen. Dengan kata lain, untuk setiap himpunan Borel
    $A_1,\ldots,A_n\subseteq\mathbb{R}$ berlaku
    \begin{align*}
       & \mathbb{P}\bigl(\xi(t_1)-\xi(t_0)\in A_1,\ldots,\xi(t_n)-\xi(t_{n-1})\in A_n\bigr) \\
       & \qquad =\prod_{k=1}^{n}\mathbb{P}\bigl(\xi(t_k)-\xi(t_{k-1})\in A_k\bigr).
    \end{align*}
  \end{definisi}

  \begin{teorema}
    Jika $\xi$ dan $\eta$ adalah variabel acak yang memiliki distribusi normal bersama
    (joint normal) pada ruang probabilitas $(\Omega,\mathcal{F},\mathbb{P})$, maka berlaku
    ekuivalensi
    \[\xi\text{ dan }\eta\text{ independen}
      \iff \mathbb{E}(\xi\eta)=\mathbb{E}(\xi)\mathbb{E}(\eta).\]
  \end{teorema}

  \begin{teorema}
    Misalkan $\xi_1,\xi_2,\ldots,\xi_n$ adalah variabel acak yang memiliki distribusi normal
    bersama (joint normal) pada ruang probabilitas $(\Omega,\mathcal{F},\mathbb{P})$. Maka
    berlaku ekuivalensi:
    \[\{\xi_1,\ldots,\xi_n\}\text{ saling independen}
      \iff \mathbb{E}(\xi_i\xi_j)=\mathbb{E}(\xi_i)\mathbb{E}(\xi_j)
      \text{ untuk setiap }i\ne j.\]
  \end{teorema}

  \begin{proposisi}
    Proses Wiener $\{W(t)\}_{t\ge 0}$ memiliki inkremen yang independen, yaitu untuk setiap
    $0=t_0<t_1<\cdots<t_n$, berlaku inkremen
    \[W(t_1)-W(t_0),\ W(t_2)-W(t_1),\ \ldots,\ W(t_n)-W(t_{n-1})\]
    saling independen.
  \end{proposisi}

  \begin{akibat}
    Jika $\{W(t)\}_{t\ge 0}$ adalah proses Wiener, maka untuk setiap
    $0=t_0<t_1<\cdots<t_n\le s<t$, berlaku bahwa $W(t)-W(s)$ independen terhadap vektor variabel
    acak
    \[\bigl(W(t_1)-W(t_0),\ W(t_2)-W(t_1),\ \ldots,\ W(t_n)-W(t_{n-1})\bigr).\]
  \end{akibat}

  \begin{akibat}
    Jika $\{W(t)\}_{t\ge 0}$ adalah proses Wiener, maka untuk setiap
    $0=t_0<t_1<\cdots<t_n\le s<t$, berlaku bahwa $W(t)-W(s)$ independen terhadap vektor variabel
    acak
    \[\bigl(W(t_1),W(t_2),\ldots,W(t_n)\bigr).\]
  \end{akibat}

  \begin{akibat}
    Jika $\{W(t)\}_{t\ge 0}$ adalah proses Wiener, maka untuk setiap $0\le s<t$, inkremen
    $W(t)-W(s)$ independen terhadap sigma-aljabar
    \[\mathcal{F}_s=\sigma\{W(r):0\le r\le s\}.\]
  \end{akibat}

  \begin{lemma}
    Proses Wiener $\{W(t)\}_{t\ge 0}$ merupakan martingale terhadap filtrasi alami
    $\{\mathcal{F}_t\}_{t\ge 0}$, yaitu
    \[\mathcal{F}_t=\sigma\{W(r):0\le r\le t\}.\]
  \end{lemma}

  \begin{lemma}
    Proses $\{|W(t)|^2-t\}_{t\ge 0}$ merupakan martingale terhadap filtrasi alami
    $\{\mathcal{F}_t\}_{t\ge 0}$, yaitu
    \[\mathcal{F}_t=\sigma\{W(r):0\le r\le t\}.\]
  \end{lemma}

  \begin{lemma}
    Proses
    \[M(t)=e^{W(t)-t/2},\qquad t\ge 0,\]
    adalah martingale terhadap filtrasi alami
    \[\mathcal{F}_t=\sigma\{W(s):0\le s\le t\}.\]
  \end{lemma}

  \begin{teorema}
    Suatu proses stokastik $W(t)$, dengan $t\ge 0$, merupakan proses Wiener jika dan hanya jika
    kondisi berikut terpenuhi:
    \begin{enumerate}[leftmargin=*, itemsep=0.5mm]
      \item hampir pasti $W(0)=0$;
      \item lintasan sampel $t\mapsto W(t)$ adalah hampir pasti kontinu;
      \item $W(t)$ memiliki inkremen yang stasioner dan saling independen;
      \item inkremen $W(t)-W(s)$ memiliki distribusi normal dengan rata-rata $0$ dan variansi
            $t-s$, untuk setiap $0\le s<t$.
    \end{enumerate}
  \end{teorema}

  \begin{lemma}
    Misalkan $X$ adalah variabel acak dengan $\mathbb{E}(X)=0$ dan
    $\mathbb{E}(X^2)=\sigma^2$. Jika untuk setiap $\lambda\in\mathbb{R}$ berlaku
    \[\mathbb{E}\bigl(e^{i\lambda X}\bigr)=\exp\!\left(-\frac{1}{2}\lambda^2\sigma^2\right),\]
    maka $X$ berdistribusi normal $\mathcal{N}(0,\sigma^2)$.
  \end{lemma}

  \begin{teorema}[Karakterisasi Martingale Levy]
    Misalkan $W(t)$, dengan $t\ge 0$, adalah proses stokastik, dan
    $\mathcal{F}_t=\sigma\bigl(W(s),s\le t\bigr)$ merupakan filtrasi yang dihasilkan oleh proses
    tersebut. Proses $\{W(t)\}_{t\ge 0}$ adalah proses Wiener jika dan hanya jika kondisi berikut
    dipenuhi:
    \begin{enumerate}[leftmargin=*, itemsep=0.5mm]
      \item hampir pasti $W(0)=0$;
      \item lintasan sampel $t\mapsto W(t)$ hampir pasti kontinu;
      \item $W(t)$ merupakan martingale terhadap filtrasi $\mathcal{F}_t$;
      \item $W(t)^2-t$ merupakan martingale terhadap filtrasi $\mathcal{F}_t$.
    \end{enumerate}
  \end{teorema}

  \begin{lemma}
    Misalkan diberikan partisi waktu
    \[0=t_0<t_1<\cdots<t_n=T,\]
    dengan $t_i=\dfrac{iT}{n}$ untuk setiap $i=0,1,\ldots,n$, dan didefinisikan inkremen
    \[\Delta_i^n W=W(t_{i+1})-W(t_i).\]
    Jika $\{W(t)\}_{t\ge 0}$ adalah proses Wiener, maka berlaku
    \[\lim_{n\to\infty}\sum_{i=0}^{n-1}(\Delta_i^n W)^2=T
      \qquad \text{dalam }L^2.\]
  \end{lemma}

  \begin{definisi}
    Variasi total dari suatu fungsi $f:[0,T]\to\mathbb{R}$ pada interval $[0,T]$ didefinisikan
    sebagai
    \[V_T(f)=\sup_{\Pi}\sum_{i=0}^{n-1}|f(t_{i+1})-f(t_i)|,\]
    dengan supremum diambil terhadap semua partisi $\Pi$ dari interval $[0,T]$, yaitu
    \[\Pi:0=t_0<t_1<\cdots<t_n=T.\]
    Besaran
    \[\|\Pi\|=\max_{0\le i\le n-1}(t_{i+1}-t_i)\]
    disebut sebagai norma partisi. Dalam pengertian limit, variasi total juga dapat dipandang
    sebagai
    \[V_T(f)=\lim_{\|\Pi\|\to 0}\sum_{i=0}^{n-1}|f(t_{i+1})-f(t_i)|,\]
    apabila limit tersebut ada. Suatu fungsi $f$ dikatakan memiliki variasi terbatas pada
    $[0,T]$ jika $V_T(f)<\infty$.
  \end{definisi}

  \begin{teorema}
    Misalkan $\{X_n\}$ adalah suatu barisan variabel acak yang didefinisikan pada ruang
    probabilitas $(\Omega,\mathcal{F},\mathbb{P})$. Jika $X_n\to X$ dalam $L^2$, maka terdapat
    suatu subsekuens $\{X_{n_k}\}$ sedemikian sehingga
    \[X_{n_k}\to X\qquad \text{hampir pasti}.\]
  \end{teorema}

  \begin{teorema}
    Variasi total dari lintasan proses Wiener $W(t)$ pada interval $[0,T]$ adalah tak hingga
    hampir pasti, yaitu
    \[V_T(W)=\infty\qquad \text{hampir pasti}.\]
  \end{teorema}

  \begin{lemma}
    Proses Wiener $\{W(t)\}_{t\ge 0}$ tidak dapat didiferensialkan pada $t=0$ hampir pasti.
  \end{lemma}

  \begin{lemma}
    Untuk setiap $t\ge 0$, proses Wiener $W(t)$ tidak dapat didiferensialkan pada titik $t$
    hampir pasti.
  \end{lemma}

  \begin{teorema}
    Dengan probabilitas $1$, lintasan proses Wiener $W(t)$ tidak terdiferensialkan untuk setiap
    $t\ge 0$.
  \end{teorema}

  \section{\small Kalkulus Ito}

  \begin{catatan}
    Dalam kasus stokastik, penjumlahan pendekatan integral
    \[\sum_{j=0}^{n-1} f(s_j)\bigl(W(t_{j+1})-W(t_j)\bigr)\]
    sangat bergantung pada pemilihan titik $s_j$. Pada integral It\^o, dipilih titik kiri
    $s_j=t_j$ agar integran hanya bergantung pada informasi masa lalu.
  \end{catatan}

  \begin{definisi}
    Suatu fungsi stokastik $f(t)$, dengan $t\ge 0$, disebut sebagai proses sederhana acak jika
    terdapat partisi berhingga
    \[0=t_0<t_1<\cdots<t_n,\]
    dan variabel acak $\eta_0,\eta_1,\ldots,\eta_{n-1}$ yang terintegralkan kuadrat, yaitu
    $\mathbb{E}(|\eta_j|^2)<\infty$, sehingga
    \[f(t)=\sum_{j=0}^{n-1}\eta_j\mathbf{1}_{[t_j,t_{j+1})}(t),\qquad t\ge 0,\]
    dengan setiap $\eta_j$ merupakan variabel acak yang $\mathcal{F}_{t_j}$-terukur.
    Himpunan semua proses sederhana acak tersebut dilambangkan dengan $M_{\text{step}}^2$.
  \end{definisi}

  \begin{definisi}
    Integral stokastik dari suatu proses sederhana acak $f\in M_{\text{step}}^2$, yang memiliki
    representasi
    \[f(t)=\sum_{j=0}^{n-1}\eta_j\mathbf{1}_{[t_j,t_{j+1})}(t),\]
    didefinisikan sebagai
    \[I^*(f)=\sum_{j=0}^{n-1}\eta_j\bigl(W(t_{j+1})-W(t_j)\bigr).\]
  \end{definisi}

  \begin{proposisi}
    Untuk setiap proses sederhana acak $f\in M_{\text{step}}^2$, integral stokastik
    $I^*(f)\in L^2$ dan berlaku isometri It\^o
    \[\mathbb{E}\bigl((I^*(f))^2\bigr)=\mathbb{E}\left(\int_0^\infty f(t)^2\,dt\right).\]
  \end{proposisi}

  \begin{lemma}
    Untuk setiap proses sederhana acak $f,g\in M_{\text{step}}^2$, berlaku
    \[\mathbb{E}\bigl(I^*(f)I^*(g)\bigr)=\mathbb{E}\left(\int_0^\infty f(t)g(t)\,dt\right).\]
  \end{lemma}

  \begin{lemma}
    Pemetaan $I^*:M_{\text{step}}^2\to L^2$ adalah linear, yaitu untuk setiap
    $f,g\in M_{\text{step}}^2$ dan $\alpha,\beta\in\mathbb{R}$ berlaku
    \[I^*(\alpha f+\beta g)=\alpha I^*(f)+\beta I^*(g).\]
  \end{lemma}

  \begin{definisi}
    Didefinisikan $M^2$ sebagai kelas proses stokastik $f(t)$, dengan $t\ge 0$, yang memenuhi:
    \begin{enumerate}[leftmargin=*, itemsep=0.5mm]
      \item
            \[\mathbb{E}\left(\int_0^\infty f(t)^2\,dt\right)<\infty.\]
      \item Terdapat barisan proses sederhana acak $\{f_n\}\subset M_{\text{step}}^2$
            sedemikian sehingga
            \[\lim_{n\to\infty}\mathbb{E}\left(\int_0^\infty (f(t)-f_n(t))^2\,dt\right)=0.\]
    \end{enumerate}
  \end{definisi}

  \begin{definisi}
    Diberikan suatu proses stokastik $f\in M^2$. Integral stokastik It\^o dari $f$ didefinisikan
    sebagai suatu variabel acak $I(f)\in L^2$ yang memenuhi
    \[\lim_{n\to\infty}\mathbb{E}\bigl(|I(f)-I^*(f_n)|^2\bigr)=0,\]
    untuk setiap barisan proses sederhana acak $\{f_n\}\subset M_{\text{step}}^2$ yang
    mendekati $f$ dalam $M^2$. Integral ini dilambangkan dengan
    \[I(f)=\int_0^\infty f(t)\,dW(t).\]
  \end{definisi}

  \begin{proposisi}
    Untuk setiap $f\in M^2$, integral stokastik $I(f)\in L^2$ ada, bersifat unik hingga hampir
    pasti, dan memenuhi
    \[\mathbb{E}\bigl((I(f))^2\bigr)=\mathbb{E}\left(\int_0^\infty f(t)^2\,dt\right).\]
  \end{proposisi}

  \begin{definisi}
    Untuk $T>0$, didefinisikan $M_T^2$ sebagai himpunan proses stokastik $f(t)$ sehingga
    $\mathbf{1}_{[0,T]}f\in M^2$. Integral stokastik It\^o pada interval $[0,T]$ untuk
    $f\in M_T^2$ didefinisikan sebagai
    \[I_T(f)=I(\mathbf{1}_{[0,T]}f)=\int_0^T f(t)\,dW(t).\]
  \end{definisi}

  \begin{lemma}
    Setiap proses sederhana acak $f\in M_{\text{step}}^2$ termasuk ke dalam $M_T^2$ untuk
    sembarang $T>0$, dan
    \[I_T(f)=\int_0^T f(s)\,dW(s)\]
    adalah martingale.
  \end{lemma}

  \begin{teorema}
    Misalkan $f(t)$, $t\ge 0$, adalah proses stokastik dengan lintasan kontinu hampir pasti,
    yang teradaptasi terhadap filtrasi $\{\mathcal{F}_t\}$. Maka:
    \begin{enumerate}[leftmargin=*, itemsep=0.5mm]
      \item $f\in M^2$, sehingga integral It\^o $I(f)$ ada, jika
            \[\mathbb{E}\left(\int_0^\infty |f(t)|^2\,dt\right)<\infty.\]
      \item $f\in M_T^2$, sehingga integral It\^o $I_T(f)$ ada, jika
            \[\mathbb{E}\left(\int_0^T |f(t)|^2\,dt\right)<\infty.\]
    \end{enumerate}
  \end{teorema}

  \begin{definisi}
    Suatu proses stokastik $f(t)$, $t\ge 0$, disebut progresif terukur jika untuk setiap
    $t\ge 0$, pemetaan
    \[(s,\omega)\mapsto f(s,\omega)\]
    merupakan fungsi yang terukur dari ruang ukur
    \[([0,t]\times\Omega,\mathcal{B}([0,t])\otimes\mathcal{F})\]
    ke $\mathbb{R}$.
  \end{definisi}

  \begin{teorema}
    Berlaku pernyataan berikut:
    \begin{enumerate}[leftmargin=*, itemsep=0.5mm]
      \item Ruang $M^2$ terdiri dari semua proses stokastik $f(t)$, $t\ge 0$, yang progresif
            terukur dan memenuhi
            \[\mathbb{E}\left(\int_0^\infty |f(t)|^2\,dt\right)<\infty.\]
      \item Ruang $M_T^2$ terdiri dari semua proses stokastik $f(t)$, $t\ge 0$, yang progresif
            terukur dan memenuhi
            \[\mathbb{E}\left(\int_0^T |f(t)|^2\,dt\right)<\infty.\]
    \end{enumerate}
  \end{teorema}

  \begin{lemma}
    Proses Wiener $W(t)$ termasuk dalam $M_T^2$ untuk setiap $T>0$.
  \end{lemma}

  \begin{teorema}
    Misalkan $T>0$, $f,g\in M_T^2$, $\alpha,\beta\in\mathbb{R}$, dan $0\le s<t\le T$.
    Maka berlaku:
    \begin{enumerate}[leftmargin=*, itemsep=0.5mm]
      \item (Linearitas)
            \[\int_0^t (\alpha f(r)+\beta g(r))\,dW(r)
              =\alpha\int_0^t f(r)\,dW(r)+\beta\int_0^t g(r)\,dW(r).\]
      \item (Isometri)
            \[\mathbb{E}\left(\left(\int_0^t f(r)\,dW(r)\right)^2\right)
              =\mathbb{E}\left(\int_0^t |f(r)|^2\,dr\right).\]
      \item (Sifat martingale)
            \[\mathbb{E}\left(\left.\int_0^t f(r)\,dW(r)\right|\mathcal{F}_s\right)
              =\int_0^s f(r)\,dW(r).\]
    \end{enumerate}
  \end{teorema}

  \begin{lemma}
    Jika $\xi,\xi_1,\xi_2,\ldots$ adalah variabel acak yang terintegrasi kuadrat dan memenuhi
    $\xi_n\to\xi$ di $L^2$ saat $n\to\infty$, maka
    \[\mathbb{E}(\xi_n\mid\mathcal{G})\to\mathbb{E}(\xi\mid\mathcal{G})
      \quad \text{di }L^2 \text{ saat }n\to\infty,\]
    untuk setiap sigma-aljabar $\mathcal{G}\subseteq\mathcal{F}$.
  \end{lemma}

  \begin{definisi}
    Misalkan $\eta(t)$ dan $\zeta(t)$ adalah proses stokastik yang didefinisikan untuk setiap
    $t\in T$, dengan $T\subseteq\mathbb{R}$. Kedua proses tersebut dikatakan saling modifikasi
    (atau versi satu sama lain) jika
    \[\mathbb{P}(\{\eta(t)=\zeta(t)\})=1,\qquad \text{untuk semua }t\in T.\]
  \end{definisi}

  \begin{teorema}
    Jika $f(s)$ adalah proses yang termasuk dalam $M_T^2$ dan didefinisikan
    \[\xi(t)=\int_0^t f(s)\,dW(s),\qquad t\ge 0,\]
    maka terdapat modifikasi teradaptasi dari $\xi(t)$ yang memiliki lintasan kontinu hampir
    pasti. Modifikasi ini unik hampir pasti.
  \end{teorema}

  \begin{definisi}
    Suatu proses stokastik $\xi(t)$ untuk $t\ge 0$ disebut proses It\^o jika memiliki lintasan
    kontinu hampir pasti, dan dapat direpresentasikan sebagai
    \[\xi(T)=\xi(0)+\int_0^T a(t)\,dt+\int_0^T b(t)\,dW(t)\qquad \text{hampir pasti},\]
    di mana $b(t)$ adalah proses yang termasuk dalam $M_T^2$ untuk semua $T\ge 0$, dan $a(t)$
    adalah proses yang teradaptasi terhadap filtrasi $\mathcal{F}_t$ sehingga
    \[\int_0^T |a(t)|\,dt<\infty\qquad \text{hampir pasti},\]
    untuk semua $T\ge 0$. Kelas semua proses teradaptasi $a(t)$ yang memenuhi syarat ini
    dilambangkan dengan $L_T^1$.
  \end{definisi}

  \begin{teorema}[Rumus Ito, Versi Sederhana]
    Misalkan $F(t,x)$ adalah fungsi bernilai real dengan turunan parsial kontinu
    $F_t(t,x)$, $F_x(t,x)$, dan $F_{xx}(t,x)$ untuk semua $t\ge 0$ dan $x\in\mathbb{R}$.
    Diasumsikan juga bahwa proses $F_x(t,W(t))$ termasuk dalam $M_T^2$ untuk semua $T\ge 0$.
    Maka $F(t,W(t))$ adalah proses It\^o yang memenuhi
    \begin{align*}
      F(T,W(T))-F(0,W(0))
       & =
      \int_0^T \left(F_t(t,W(t))+\frac{1}{2}F_{xx}(t,W(t))\right)\,dt \\
       & \quad +
      \int_0^T F_x(t,W(t))\,dW(t)
      \qquad \text{hampir pasti}.
    \end{align*}
    Dalam notasi diferensial,
    \[dF(t,W(t))
      =\left(F_t(t,W(t))+\frac{1}{2}F_{xx}(t,W(t))\right)dt+F_x(t,W(t))\,dW(t).\]
  \end{teorema}

  \begin{teorema}[Rumus Ito, Kasus Umum]
    Misalkan $\xi(t)$ adalah proses It\^o yang memenuhi
    \[d\xi(t)=a(t)\,dt+b(t)\,dW(t).\]
    Andaikan $F(t,x)$ adalah fungsi bernilai real dengan turunan parsial kontinu
    $F_t(t,x)$, $F_x(t,x)$, dan $F_{xx}(t,x)$ untuk semua $t\ge 0$ dan $x\in\mathbb{R}$.
    Jika proses $b(t)F_x(t,\xi(t))$ termasuk dalam $M_T^2$ untuk setiap $T\ge 0$, maka
    $F(t,\xi(t))$ merupakan proses It\^o dan memenuhi
    \begin{align*}
      dF(t,\xi(t))
       & =
      \left(F_t(t,\xi(t))+F_x(t,\xi(t))a(t)+\frac{1}{2}F_{xx}(t,\xi(t))b(t)^2\right)\,dt \\
       & \quad +F_x(t,\xi(t))b(t)\,dW(t).
    \end{align*}
  \end{teorema}

  \begin{teorema}
    Misalkan
    \[dX_t=\mu(t,X_t)\,dt+\sigma(t,X_t)\,dW_t,\]
    dengan:
    \begin{itemize}[leftmargin=*, itemsep=0.5mm]
      \item $X_t=(X_t^1,\ldots,X_t^n)^{\mathsf T}\in\mathbb{R}^n$ adalah proses It\^o
            berdimensi $n$,
      \item $\mu(t,X_t)=(\mu_1(t,X_t),\ldots,\mu_n(t,X_t))^{\mathsf T}\in\mathbb{R}^n$ adalah
            vektor drift,
      \item $\sigma(t,X_t)=[\sigma_{ij}(t,X_t)]_{n\times m}\in\mathbb{R}^{n\times m}$ adalah
            matriks volatilitas,
      \item $W_t=(W_t^1,\ldots,W_t^m)^{\mathsf T}\in\mathbb{R}^m$ adalah proses Wiener
            berdimensi $m$.
    \end{itemize}
    Jika $f\in C^{1,2}([0,\infty)\times\mathbb{R}^n)$, maka berlaku
    \begin{align*}
      df(t,X_t)
       & =
      \frac{\partial f}{\partial t}(t,X_t)\,dt+(\nabla f(t,X_t))^{\mathsf T}\mu(t,X_t)\,dt \\
       & \quad +
      (\nabla f(t,X_t))^{\mathsf T}\sigma(t,X_t)\,dW_t
      +\frac{1}{2}\operatorname{Tr}\!\bigl(\sigma(t,X_t)\sigma(t,X_t)^{\mathsf T}D^2f(t,X_t)\bigr)\,dt.
    \end{align*}
  \end{teorema}

  \begin{definisi}
    Suatu proses It\^o $\xi(t)$, $t\ge 0$, disebut solusi dari masalah nilai awal
    \[d\xi(t)=f(\xi(t))\,dt+g(\xi(t))\,dW(t),\qquad \xi(0)=\xi_{\mathrm{init}},\]
    jika $\xi_{\mathrm{init}}$ adalah variabel acak $\mathcal{F}_0$-terukur,
    proses $f(\xi(t))$ dan $g(\xi(t))$ masing-masing termasuk dalam $L_T^1$ dan $M_T^2$,
    serta
    \[\xi(T)=\xi_{\mathrm{init}}+\int_0^T f(\xi(t))\,dt+\int_0^T g(\xi(t))\,dW(t)
      \qquad \text{hampir pasti}\]
    untuk semua $T\ge 0$.
  \end{definisi}

  \begin{teorema}
    Misalkan $f$ dan $g$ adalah fungsi Lipschitz kontinu dari $\mathbb{R}$ ke $\mathbb{R}$,
    yaitu terdapat konstanta $C>0$ sehingga untuk semua $x,y\in\mathbb{R}$ berlaku
    \[|f(x)-f(y)|\le C|x-y|,\qquad |g(x)-g(y)|\le C|x-y|.\]
    Selain itu, misalkan $\xi_{\mathrm{init}}$ adalah variabel acak $\mathcal{F}_0$-terukur
    yang terintegral kuadrat. Maka masalah nilai awal
    \[d\xi(t)=f(\xi(t))\,dt+g(\xi(t))\,dW(t),\qquad \xi(0)=\xi_{\mathrm{init}},\]
    memiliki solusi $\xi(t)$, $t\ge 0$, dalam kelas proses It\^o. Solusi ini unik dalam arti
    bahwa jika $\eta(t)$, $t\ge 0$, adalah solusi lain, maka
    \[\mathbb{P}(\xi(t)=\eta(t),\ \forall t\ge 0)=1.\]
    Lebih lanjut, diperoleh bahwa $\xi(t)\in L^2$ untuk semua $t\in[0,T]$, yaitu
    \[\mathbb{E}|\xi(t)|^2\le
      e^{C^2t(3t+6)}\Bigl(3\mathbb{E}|\xi_{\mathrm{init}}|^2+6t(|f(0)|^2+|g(0)|^2)\Bigr).\]
  \end{teorema}

  \begin{teorema}[Titik Tetap Banach]
    Misalkan $(X,d)$ adalah ruang metrik lengkap dan $T:X\to X$ adalah suatu pemetaan kontraksi,
    yaitu terdapat konstanta $0\le \alpha<1$ sehingga
    \[d(T(x),T(y))\le \alpha\, d(x,y)\qquad \text{untuk semua }x,y\in X.\]
    Maka:
    \begin{enumerate}[leftmargin=*, itemsep=0.5mm]
      \item $T$ memiliki tepat satu titik tetap $x^*\in X$, yaitu $T(x^*)=x^*$.
      \item Untuk setiap $x_0\in X$, barisan yang didefinisikan dengan iterasi
            \[x_{n+1}=T(x_n)\]
            konvergen ke $x^*$.
      \item Selain itu, terdapat taksiran konvergensi
            \[d(x_n,x^*)\le \frac{\alpha^n}{1-\alpha}d(x_1,x_0).\]
    \end{enumerate}
  \end{teorema}

  \begin{teorema}
    Misalkan $\xi(t)$ adalah solusi eksak dari persamaan diferensial stokastik
    \[d\xi(t)=f(\xi(t))\,dt+g(\xi(t))\,dW(t),\qquad \xi(0)=\xi_{\mathrm{init}},\]
    dan $\xi_k^n$ adalah solusi numerik dari metode Euler-Maruyama dengan partisi waktu.
    Jika $f$ dan $g$ memenuhi kondisi Lipschitz global dan linear growth, maka terdapat
    konstanta $K>0$ yang bergantung pada $C,f(0),g(0),\xi_{\mathrm{init}}$, dan $T$, sehingga
    \[\|\xi(t_k^n)-\xi_k^n\|_{L^2}
      \le K\left(\sup_{i=0,1,\ldots,k-1}\Delta_i^n t\right)^{1/2}.\]
  \end{teorema}
\end{multicols}
\end{document}
